\chapter{半双线性型与二次型}

\section{半双线性型}

\subsection{双线性映射。}

本小节中，A 与 B 表示两个环，E 表示左 A-模，F 表示右 B-模，G 表示 (A, B)-双模，即一个交换群，配备左 A-模结构与右 B-模结构，使得 \((ag)b = a(gb)\) 对一切 \(a \in A\) 、\(b \in B\) 、\(g \in G\) 成立。

定义 1. --- 称一个映射 \(\Phi\) 为从积 \(\mathbf{E} \times \mathbf{F}\) 到 \(\mathbf{G}\) 中的双线性映射，如果它满足下列条件：

\begin{enumerate}
\def\labelenumi{(\arabic{enumi})}
\tightlist
\item
  \(\Phi(x + x', y) = \Phi(x, y) + \Phi(x', y)\)
\end{enumerate}

对一切 \(x \in \mathbf{E}\) 、\(x' \in \mathbf{E}\) 、\(y \in \mathbf{F}\) ；

\begin{enumerate}
\def\labelenumi{(\arabic{enumi})}
\setcounter{enumi}{1}
\tightlist
\item
  \(\Phi(x, y + y') = \Phi(x, y) + \Phi(x, y')\)
\end{enumerate}

对一切 \(x \in \mathbf{E}, y \in \mathbf{F}, y' \in \mathbf{F}\) ；

\begin{enumerate}
\def\labelenumi{(\arabic{enumi})}
\setcounter{enumi}{2}
\item
  \(\Phi(ax, y) = a\Phi(x, y)\) 对一切 \(a \in \mathbf{A}\) 、\(x \in \mathbf{E}\) 、\(y \in \mathbf{F}\) ；
\item
  \(\Phi(x, yb) = \Phi(x, y)b\) 对一切 \(x \in \mathbf{E}, y \in \mathbf{F}, b \in \mathbf{B}\) 。
\end{enumerate}

张量积 \(E \otimes_{Z} F\) 上典范地带有由 \(a(x \otimes y)b = ax \otimes yb\) 刻画的一个 (A, B)-双模结构（第 III 章，第 2 版，附录 II，第 3 小节），而双线性映射 \(\Phi\) （从 \(E \times F\) 到 G 中）的数据，等价于映射 \(\Psi\) （从 \(E \otimes_{Z} F\) 到 G 中，对于 (A, B)-双模结构而言是同态，且满足 \(\Psi(x \otimes y) = \Phi(x, y)\) ，对一切 \(x \in E\) 与 \(y \in F\)）的数据。

定义 1 加于 \(\Phi\) 的条件表明，偏映射 \(d_{\Phi}(y): x \to \Phi(x, y)\) 与 \(s_{\Phi}(x): y \to \Phi(x, y)\) 分别是 E 到 G 的 A-线性映射与 F 到 G 的 B-线性映射。给交换群 \(\mathcal{L}_{\mathbf{A}}(\mathbf{E}, \mathbf{G})\) （相应地 \(\mathcal{L}_{\mathbf{B}}(\mathbf{F}, \mathbf{G})\) ）配上由 \(ub(x) = u(x). b\) （ \(u \in \mathcal{L}_{\mathbf{A}}(\mathbf{E}, \mathbf{G}), x \in \mathbf{E}, b \in \mathbf{B}\) ）（相应地 \(av(y) = a. v(y)\) （ \(a \in \mathbf{A}, v \in \mathcal{L}_{\mathbf{B}}(\mathbf{F}, \mathbf{G}), y \in \mathbf{F}\) ））所定义的右 B-模（相应地左 A-模）结构。于是，条件 (1) 至 (4) 分别等价于：

\[
s _ {\Phi} (x + x ^ {\prime}) = s _ {\Phi} (x) + s _ {\Phi} (x ^ {\prime})\tag{1'}
\]

\[
d _ {\Phi} (y + y ^ {\prime}) = d _ {\Phi} (y) + d _ {\Phi} (y ^ {\prime})\tag{2'}
\]

\[
s _ {\Phi} (a x) = a. s _ {\Phi} (x)\tag{3'}
\]

\[
d _ {\Phi} (y b) = d _ {\Phi} (y). b,\tag{4'}
\]

对 E 中的一切 \(x, x'\) 、F 中的一切 \(y, y'\) 、\(a \in A\) 、\(b \in B\) ；换言之，映射 \(d_{\Phi}\) （从 F 到 \(\mathcal{L}_{A}(E, G)\) 中）是 B-线性的，而映射 \(s_{\Phi}\) （从 E 到 \(\mathcal{L}_{B}(F, G)\) 中）是 A-线性的。按定义，我们有

\begin{enumerate}
\def\labelenumi{(\arabic{enumi})}
\setcounter{enumi}{4}
\tightlist
\item
  \(\Phi(x, y) = d_{\Phi}(y)(x) = s_{\Phi}(x)(y)\) 对一切 \(x \in \mathbf{E}, y \in \mathbf{F}\) 。
\end{enumerate}

定义 2. --- 给定 E × F 到 G 的双线性映射 \(\Phi\)，映射 \(d_{\Phi}\) （从 F 到 \(\mathfrak{L}_{\mathrm{A}}(\mathrm{E}, \mathrm{G})\) 中；相应地，映射 \(s_{\Phi}\) 从 E 到 \(\mathfrak{L}_{\mathrm{B}}(\mathrm{F}, \mathrm{G})\) 中）由 (5) 刻画，称为 \(\Phi\) 的右相伴线性映射（相应地左相伴线性映射）。

反之，B-线性映射 \(d\) （从 F 到 \(\mathfrak{L}_{\mathbf{A}}(\mathbf{E},\mathbf{G})\) 中；相应地，A-线性映射 \(s\) 从 E 到 \(\mathfrak{L}_{\mathbf{B}}(\mathbf{F},\mathbf{G})\) 中）的数据，借助公式

\[
\Phi (x, y) = d (y) (x) \quad (\text { resp. } \Phi (x, y) = s (x) (y))
\]

唯一确定一个双线性映射 \(\Phi\)，它映 \(\mathbf{E} \times \mathbf{F}\) 到 \(\mathbf{G}\) 中，且 \(d\) （相应地 \(s\)）是它的右相伴线性映射（相应地左相伴线性映射）。

定义 3. --- 称 E × F 到 G 的双线性映射 \(\Phi\) 为右退化的（相应地左退化的），如果存在 F 的一个非零元素 \(y_0\) （相应地 E 的一个非零元素 \(x_0\)），使得 \(\Phi(x, y_0) = 0\) 对一切 \(x \in \mathbf{E}\) 成立（相应地 \(\Phi(x_0, y) = 0\) 对一切 \(y \in \mathbf{F}\) 成立）。称 \(\Phi\) 为退化的，如果它是右退化的或左退化的。

\(\Phi\) 右非退化（相应地左非退化）必须且只需 \(\Phi\) 的右相伴线性映射（相应地左相伴线性映射）是单射；因此，说 \(\Phi\) 非退化，就是说相伴线性映射 \(d_{\Phi}\) 与 \(s_{\Phi}\) 都是单射。

设 \((e_i)_{i\in I}\) 与 \((f_k)_{k\in K}\) 是 E 与 F 的元素的两个族，并设 \((a_i)_{i\in I}\) 与 \((b_k)_{k\in K}\) 是 A 与 B 的元素的两个族，其中除有限个外均为零。由等式 (1) 至 (4)，对非零系数的个数作归纳，可得

\[
\Phi (\sum_ {i} a _ {i} e _ {i}, \sum_ {k} f _ {k} b _ {k}) = \sum_ {i, k} a _ {i} \Phi (e _ {i}, f _ {k}) b _ {k}.\tag{6}
\]

若 \((e_{i})\) 与 \((f_{k})\) 是模 E 与 F 的生成系，则 \(\Phi\) 由元素 \(g_{ik} = \Phi(e_{i}, f_{k})\) 完全确定。若 \((e_{i})\) 与 \((f_{k})\) 是 E 与 F 的基，且给定了 G 的元素 \(g_{ik}\) \((i \in I, k \in K)\)，则公式

\[
\Phi (\sum_ {i} a _ {i} e _ {i}, \sum_ {k} f _ {k} b _ {k}) = \sum_ {i, k} a _ {i} g _ {i k} b _ {k}\tag{6'}
\]

定义了 E × F 到 G 的一个映射，它是双线性的，且满足 \(\Phi(e_{i}, f_{k}) = g_{ik}\) 。当 \((e_{i})\) 与 \((f_{k})\) 为有限基时，称 \((\Phi(e_{i}, f_{k}))\) 为 \(\Phi\) 关于这些基的矩阵。

E × F 到 G 中的双线性映射显然构成 E × F 到 G 的映射的加法群的一个子群。另一方面，设 a（相应地 b）为 A（相应地 B）的中心的一个元素；则映射 \(a\Phi b\) （从 E × F 到 G 中，由 \((a\Phi b)(x, y) = a \cdot \Phi(x, y) \cdot b\) 定义）是双线性的。于是，E × F 到 G 中的双线性映射的集合便带有 A 与 B 的中心上的一个双模结构。

设 E′（相应地 F′）为左 A-模（相应地右 B-模），u（相应地 v）为 E 到 E′（相应地 F 到 F′）的同态，\(\Phi'\) 为 \(E' \times F'\) 到 G 中的双线性映射。\(\Phi'\) （关于 u 与 v）的原像是 E × F 到 G 的双线性映射 \(\Phi\)，它定义为

\[
\Phi (x, y) = \Phi^ {\prime} (u (x), v (y)) \quad (x \in \mathrm{E}, y \in \mathrm{F}).
\]

容易验证，

\[
d _ {\Phi} (y) = d _ {\Phi^ {\prime}} (\nu (y)) \circ u \quad \text { and } \quad s _ {\Phi} (x) = s _ {\Phi^ {\prime}} (u (x)) \circ \nu
\]

对一切 \(x \in \mathbf{E}, y \in \mathbf{F}\) 。

设 \(\Phi\) 为 \(\mathbf{E} \times \mathbf{F}\) 到 \(\mathbf{G}\) 中的双线性映射，\(h\) 为 \(\mathbf{G}\) 到另一个 (A, B)-双模 \(\mathbf{G}'\) 中的一个（对于 (A, B)-双模结构的）同态。则 \(h \circ \Phi\) 是 \(\mathbf{E} \times \mathbf{F}\) 到 \(\mathbf{G}'\) 中的双线性映射。

\subsection{半双线性映射。}

本节中，除非另有明确说明，我们用 A 与 B 表示两个环，用 E 表示左 A-模，用 F 表示左 B-模；我们用 \(b \to b^{\mathrm{J}} (b \in \mathrm{B})\) 表示 B 的一个反自同构，即 B 到其自身上的一个双射，满足 \((b + c)^{\mathrm{J}} = b^{\mathrm{J}} + c^{\mathrm{J}}\) 与 \((bc)^{\mathrm{J}} = c^{\mathrm{J}}b^{\mathrm{J}}\) （对 B 中的一切 \(b, c\)）；我们用 \(\mathbf{J}'\) 表示 \(\mathbf{J}^{-1}\) 。我们用 G 表示一个 \((\mathbf{A}, \mathbf{B})\) -双模（第 1 小节）。

定义 4. --- 称映射 \(\Phi\) （从 E \(\times\) F 到 G 中）为关于 J 的右半双线性映射，如果它满足条件 (1)、(2)、(3)（定义 1，第 1 小节）以及

\begin{enumerate}
\def\labelenumi{(\arabic{enumi})}
\setcounter{enumi}{6}
\tightlist
\item
  \(\Phi(x, by) = \Phi(x, y). b^J\) 对一切 \(x \in \mathbf{E}, y \in \mathbf{F}\) 与 \(b \in \mathbf{B}\) 。
\end{enumerate}

若 J 是恒同映射（这要求 B 是交换的），我们就重新得到双线性映射的概念。

设 \((e_i)_{i\in I}\) 与 \((f_k)_{k\in K}\) 是 E 与 F 的元素的两个族，并设 \((a_i)_{i\in I}\) 与 \((b_k)_{k\in K}\) 是 A 与 B 的元素，其中除有限个外均为零。我们有

\[
\Phi (\sum_ {i} a _ {i} e _ {i}, \sum_ {k} b _ {k} f _ {k}) = \sum_ {i, k} a _ {i} \Phi (e _ {i}, f _ {k}) b _ {k} ^ {\mathrm{J}}.\tag{8}
\]

与双线性映射的情形一样，元素 \(\Phi(e_i, f_k)\) 以唯一的方式确定 \(\Phi\) （当 \((e_i)\) 与 \((f_k)\) 是生成系时），而当 \((e_i)\) 与 \((f_k)\) 是 E 与 F 的基时可任意选取；当 \((e_i)\) 与 \((f_k)\) 是有限基时，称 \((\Phi(e_i, f_k))\) 为 \(\Phi\) 关于这些基的矩阵。

与双线性映射一样，在 E × F 到 G 中的（关于 J 的）右半双线性映射的集合上，可定义 A 与 B 的中心上的一个双模结构。半双线性映射的原像概念，用与双线性映射相同的公式来定义。此外我们将看到，半双线性映射的研究可以化为双线性映射的研究。

定义 5. --- 设 B 是一个环，F 是左（相应地右）B-模，J 是 B 的一个反自同构。用 F \(^{j}\) 表示具有与 F 相同的底加法群、且外合成律为 \((b, y) \to b^{J'}y\) （相应地 \((b, y) \to yb^{J'}\) ）\((b \in B, y \in F, J' = J^{-1})\) 的右（相应地左）B-模。

按定义 5 的记号，从 \(\mathbf{F}^{\mathrm{j}}\) 到一个右（相应地左）B-模 H 中的线性映射，因而就等同于 F 到 H 的一个 Z-线性映射 \(u\) ，它满足

\[
u (b y) = u (y) b ^ {J} \quad (\text { resp. } u (y b) = b ^ {J} u (y)) \quad (b \in B, y \in F).
\]

映射 \(u\) 是 F 到 H 的、关于 J 的半线性映射（第 II 章，附录 I，第 1 小节），如果把 J 看作 B 的反环 B \(^{0}\) 到 B 上的一个同构，并把 F 看作一个右（相应地左）B \(^{0}\) -模。

类似地，当 F 是左 B-模时，E × F 到 G 中的（关于 J 的）右半双线性映射 \(\Phi\) 等同于 E × F \(^{J}\) 到 G 中的一个双线性映射；若后者右退化（相应地左退化、非退化），则称 \(\Phi\) 右退化（相应地左退化、非退化）。

注记. --- 设 A 与 B 为两个环，\(J_{1}\) 为 A 的一个反自同构，M 为右 A-模，N 为右 B-模，G 为 (A, B)-双模。称映射 \(\Phi\) （从 M \(\times\) N 到 G 中）为关于 J \(_{1}\) 的左半双线性映射，如果它是 Z-双线性的且满足

\[
\Phi (x a, y b) = a ^ {\mathrm{J} _ {1}} \Phi (x, y) b \quad (x \in \mathrm{M}, y \in \mathrm{N}, a \in \mathrm{A}, b \in \mathrm{B}).\tag{9}
\]

这样的映射等同于 \(M^{J_{1}} \times N\) 到 G 中的一个双线性映射。我们常把为右半双线性映射给出的定义与性质向左半双线性映射的转述留给读者；当我们说到半双线性映射（不加指明）时，它总是右半双线性映射。

\subsection{正交性。双线性或半双线性映射的直和。}

本节中，A 与 B 表示环，E 表示左 A-模，F 表示右（相应地左）B-模，G 表示 (A, B)-双模，Φ 表示 E × F 到 G 的双线性（相应地，关于 B 的一个给定反自同构 J 的半双线性）映射。

定义 6. --- 称两个元素 \(x \in E\) 与 \(y \in F\) 关于 \(\Phi\) 正交，如果 \(\Phi(x, y) = 0\) 。称两个子集 \(E' \subset E\) 与 \(F' \subset F\) 正交，如果对任意 \(x \in E'\) 与 \(y \in F'\) ，\(x\) 与 \(y\) 正交。E（相应地 F）中与 F 的给定子模 N（相应地 E 的子模 M）正交的元素所成之集是 E（相应地 F）的一个子模，称为与 N（相应地 M）完全正交（或简称正交）的子模，记作 \(N^0\) （相应地 \(M^0\) ）。

设 H 与 H′ 是 E 或 F 的两个子模。有 \(\mathrm{H}\subset(\mathrm{H}^{0})^{0}\) （记作 \(H^{00}\) ）；若 \(H\subset H'\) ，则 \(H^{\prime0}\supset H^{0}\) 。由此得 \(H^{0}\supset(H^{00})^{0}\) 且 \(H^{0}\subset(H^{0})^{00}\) ；令

\[
\mathrm{H} ^ {0 0 0} = (\mathrm{H} ^ {0 0}) ^ {0} = (\mathrm{H} ^ {0}) ^ {0 0} = ((\mathrm{H} ^ {0}) ^ {0}) ^ {0},
\]

便有 \(\mathbf{H}^0 = \mathbf{H}^{000}\) 。

映射 \(\Phi\) 退化（第 1 小节，定义 3）必须且只需两个子模 E \(^{0}\) 、F \(^{0}\) 中至少有一个 \(\neq \{0\}\) 。显然，\(\Phi(x, y)\) 在给 \(x\) （相应地 \(y\) ）加上

F \(^{0}\) （相应地 E \(^{0}\) ）中的一个元素时不改变，因而 Φ 通过过渡到商，在 (E/F \(^{0}\) ) × (F/E \(^{0}\) ) 上定义了一个双线性（或半双线性）映射；它显然非退化；称之为 Φ 的相伴非退化双线性（或半双线性）映射。

设 \((\mathrm{E}_{i})_{i\in\mathrm{I}}\) 为左 A-模的一个族，\((\mathrm{F}_{i})_{i\in\mathrm{I}}\) 为右 B-模（相应地左 B-模）的一个族，\(\Phi_{i}\) 为 \(E_{i}\times F_{i}\) 到 G 中的一个双线性（相应地关于 J 的右半双线性）映射。用 E（相应地 F）表示诸 \(E_{i}\) （相应地 \(F_{i}\) ）的直和模。我们立即看到，映射 \(\Phi\) （从 \(E\times F\) 到 G 中）由

\[
\Phi ((x _ {i}), (y _ {i})) = \sum_ {i} \Phi_ {i} (x _ {i}, y _ {i}) \quad (x _ {i} \in \mathrm{E} _ {i}, y _ {i} \in \mathrm{F} _ {i})\tag{10}
\]

（这个和有意义，因为除有限项外其各项均为零）定义，它是双线性的（相应地关于 J 是右半双线性的）。称它为诸映射 \(\Phi_{i}\) 的直和。显然，\(\mathbf{E}_{i}\) 与 \(\mathbf{F}_{j}\) 关于 \(\Phi\) 正交（\(i \neq j\) ）。反之，设 \(\Phi\) 为 \(\mathbf{E} \times \mathbf{F}\) 到 G 中的一个双线性或半双线性映射，并设 E 是子模 \((\mathbf{E}_{i})_{i \in I}\) 的直和，F 是子模 \((\mathbf{F}_{i})_{i \in I}\) 的直和，使得 \(\mathbf{E}_{i}\) 与 \(\mathbf{F}_{j}\) 正交（对 \(i \neq j\)）；则 \(\Phi\) 是它在积 \(\mathbf{E}_{i} \times \mathbf{F}_{i} (i \in I)\) 上的限制的直和。

\(\Phi\) 非退化必须且只需诸 \(\Phi_{i}\) 都非退化；在这些条件下，与 \(\mathbf{E}_{i}\) 正交的子模是 \(\sum_{j\neq i}\mathbf{F}_{j}\) 。

\subsection{更换基环。}

本节中，我们用 A、B、A′、B′ 表示四个环，用 h 与 h′ 分别表示 A 到 A′ 与 B 到 B′ 的同态，用 G 表示 (A, B)-双模，用 G′ 表示 (A′, B′)-双模，并用 u 表示 G 的底加法群到 G′ 的底加法群的一个同态，它满足

\[
u (a g b) = h (a) u (g) h ^ {\prime} (b) \quad (a \in \mathrm{A}, g \in \mathrm{G}, b \in \mathrm{B}).\tag{11}
\]

设 E（相应地 F）为左 A-模（相应地右 B-模）。回顾（第 III 章，第 2 版，附录 II，第 10 小节）：若把 A′（相应地 B′）看作右 A-模（相应地左 B-模），则张量积 \(E' = A' \otimes_{A} E\) （相应地 \(F' = F \otimes_{B} B'\) ）带有由下式定义的左 A′-模（相应地右 B′-模）结构

\[
\begin{array}{c c} a _ {1} ^ {\prime} (a ^ {\prime} \otimes x) = (a _ {1} ^ {\prime} a ^ {\prime}) \otimes x & (a ^ {\prime}, a _ {1} ^ {\prime} \in A ^ {\prime}, x \in E) \\ (\text { resp. } (y \otimes b ^ {\prime}) b _ {1} ^ {\prime} = y \otimes (b ^ {\prime} b _ {1} ^ {\prime}) & (b ^ {\prime}, b _ {1} ^ {\prime} \in B ^ {\prime}, y \in F)). \end{array}\tag{12}
\]

命题 1. --- 设 E 为左 A-模，F 为右 B-模；令 \(\mathbf{E}^{\prime} = \mathbf{A}^{\prime}\otimes_{\mathbf{A}}\mathbf{E}\) 与 \(\mathbf{F}^{\prime} = \mathbf{F}\otimes_{\mathbf{B}}\mathbf{B}^{\prime}\) 。对每一个双线性映射 \(\Phi\) （从 \(\mathbf{E}\times \mathbf{F}\) 到 G 中），存在且仅存在一个双线性映射 \(\Phi^{\prime}\) （从 \(\mathbf{E}^{\prime}\times \mathbf{F}^{\prime}\) 到 G′ 中），使得有

\[
\Phi^ {\prime} (a ^ {\prime} \otimes x, y \otimes b ^ {\prime}) = a ^ {\prime}. u (\Phi (x, y)). b ^ {\prime}\tag{13}
\]

对一切 \(a' \in \mathrm{A}'\) 、\(b' \in \mathrm{B}'\) 、\(x \in \mathrm{E}\) 、\(y \in \mathrm{F}\) 。

\(\Phi'\) 的唯一性由元素 \(a' \otimes x\) 与 \(y \otimes b'\) 分别生成 E′ 与 F′ 这一事实得出。为证其存在性，考虑映射

\[
m: (a ^ {\prime}, x, y, b ^ {\prime}) \rightarrow a ^ {\prime}. u (\Phi (x, y)). b ^ {\prime}
\]

它把 A′ × E × F × B′ 映到 G 中；它显然是 Z-重线性的，且满足

\[
\begin{array}{r} m (a ^ {\prime}, a x, y, b ^ {\prime}) = m (a ^ {\prime} h (a), x, y, b ^ {\prime}) \\ m (a ^ {\prime}, x, y b, b ^ {\prime}) = m (a ^ {\prime}, x, y, h ^ {\prime} (b) b ^ {\prime}) \\ (a \in \mathrm{A}, b \in \mathrm{B}, a ^ {\prime} \in \mathrm{A} ^ {\prime}, b ^ {\prime} \in \mathrm{B} ^ {\prime}, x \in \mathrm{E}, y \in \mathrm{F}). \end{array}
\]

因此存在 E′ × F′ 到 G′ 中的、满足 (13) 的一个 Z-双线性映射 \(\Phi'\) （第 III 章，第 2 版，附录 II，第 1 小节，命题 2）。这一关系以及由 (12) 给出的 E′ 与 F′ 的模结构定义表明 \(\Phi'\) 是双线性的，证毕。

在命题 1 的假设与记号下，我们现在来研究与 \(\Phi\) 及 \(\Phi'\) 相伴的线性映射（第 1 小节，定义 2）。为此，我们先定义 \(\mathfrak{L}_{A}(\mathrm{E}, \mathrm{G})\) 到 \(\mathfrak{L}_{A'}(\mathrm{E}', \mathrm{G}')\) 中的一个典范同态。对每个 \(\nu \in \mathfrak{L}_{\Lambda}(\mathrm{E}, \mathrm{G})\) ，映射 \((a', x) \to a'. u(\nu(x))\) （从 \(\mathrm{A}' \times \mathrm{E}\) 到 \(\mathrm{G}'\) 中）是 Z-双线性的，且鉴于 (11)，它把 \((a'h(a), x)\) 与 \((a', ax)\) （ \(a \in \mathrm{A}\) ）映到 \(\mathrm{G}'\) 的同一元素上；因而它定义了（第 III 章，第 2 版，

附录 II，第 1 与 10 小节）一个映射 \(k(\nu)\) （从 \(\mathbf{E}' = \mathbf{A}' \otimes_{\mathbf{A}} \mathbf{E}\) 到 \(\mathbf{G}'\) 中），使得 \(k(\nu)(a' \otimes x) = a'. u(\nu(x))\) ，且由 (12)，它是 A′-线性的。此外，由 (12) 立即推出，映射 \(\nu \to k(\nu)\) （从 \(\mathfrak{L}_{\mathbf{A}}(\mathbf{E}, \mathbf{G})\) 到 \(\mathfrak{L}_{\mathbf{A}'}(\mathbf{E}', \mathbf{G}')\) 中）满足 \(k(\nu b) = k(\nu) h'(b)\) （对一切 \(b \in B\)）。用 \(i\) 表示典范映射 \(y \to y \otimes 1\) （从 F 到 \(\mathbf{F}'\) 中）。则图表

\[
\begin{array}{c c c} \mathrm{F} & \xrightarrow {d _ {\Phi}} & \mathcal {L} _ {A} (\mathrm{E}, \mathrm{G}) \\ \downarrow i & & \downarrow k \\ \mathrm{F} ^ {\prime} & \xrightarrow {d _ {\Phi^ {\prime}}} & \mathcal {L} _ {A^{\prime}} (\mathrm{E} ^ {\prime}, \mathrm{G} ^ {\prime}) \end{array}\tag{14}
\]

（其中 \(d_{\Phi}\) 与 \(d_{\Phi'}\) 表示 \(\Phi\) 与 \(\Phi'\) 的右相伴线性映射）是交换的。事实上，对 \(x \in \mathrm{E}, y \in \mathrm{F}\) 与 \(a' \in \mathrm{A}'\) ，有 \(d_{\Phi'}(i(y))(a' \otimes x) = \Phi'(a' \otimes x, y \otimes 1) = a'. u(\Phi(x, y)) = a'. u(d_{\Phi}(y)(x))\) ，也就是说 \(d_{\Phi'}(i(y))(a' \otimes x) = k(d_{\Phi}(y))(a' \otimes x)\) 。线性映射 \(s_{\Phi}\) 与 \(s_{\Phi'}\) （它们左相伴于 \(\Phi\) 与 \(\Phi'\) ）有类似的交换关系。

命题 2. --- 设 B 与 B′ 分别配备了反自同构 J 与 I，使得

\[
h ^ {\prime} (b ^ {J}) = h ^ {\prime} (b) ^ {I} \quad \text { for all } b \in \mathrm{B}.\tag{15}
\]

设 E 为左 A-模，F 为左 B-模；令 \(E' = A' \otimes_{A} E\) 与 \(F' = B' \otimes_{B} F\) 。对每一个（关于 J 的）半双线性映射 \(\Phi\) （从 \(E \times F\) 到 G 中），存在且仅存在一个（关于 I 的）半双线性映射 \(\Phi'\) （从 \(E' \times F'\) 到 \(G'\) 中），使得

\[
\Phi^ {\prime} (a ^ {\prime} \otimes x, b ^ {\prime} \otimes y) = a ^ {\prime}. u (\Phi (x, y)). b ^ {\prime I}\tag{16}
\]

对一切 \(a' \in \mathbf{A}'\) 、\(b' \in \mathbf{B}'\) 、\(x \in \mathbf{E}\) 、\(y \in \mathbf{F}\) 。

\(\Phi'\) 的唯一性由张量积 \(a' \otimes x\) 与 \(b' \otimes y\) 分别生成 E′ 与 F′ 这一事实得出。为证其存在性，考虑映射

\[
m: (a, x, b ^ {\prime}, y) \rightarrow a ^ {\prime}. u (\Phi (x, y)). b ^ {\prime I}
\]

它把 \(\mathbf{A}' \times \mathbf{E} \times \mathbf{B}' \times \mathbf{F}\) 映到 \(\mathbf{G}\) 中。它显然是 \(\mathbf{Z}\)-重线性的，且把 (11) 与 (15) 考虑在内，满足 \(m(a', ax, b', y) = m(a'h(a), x, b', y)\) 与 \(m(a', x, b', by) = a'. u(\Phi(x, y)) \cdot h'(b)^{I} b'^{I} = m(a', x, b'h'(b), y)\) （ \(a \in \mathbf{A}, b \in \mathbf{B}, a' \in \mathbf{A}'\) ， \(b' \in \mathbf{B}'\) ， \(x \in \mathbf{E}\) ， \(y \in \mathbf{F}\) ）。因此存在一个 \(\mathbf{Z}\)-双线性映射 \(\Phi'\) （从 \(\mathbf{E}' \times \mathbf{F}'\) 到 \(\mathbf{G}'\) 中）

它满足 (16) （第 III 章，第 2 版，附录 II，第 1 小节，命题 2）。这一关系连同由 (12) 给出的 E′ 与 F′ 的模结构定义，并鉴于 (15)，表明 Φ′ 是关于 I 半双线性的，证毕。

\((\mathbf{A}^{\prime}, \mathbf{B}^{\prime})\) -双模 \(G^{\prime}\) （配备了满足 (11) 的、G 到 \(G^{\prime}\) 中的 Z-线性映射 u）的最重要例子如下：

\begin{enumerate}
\def\labelenumi{\arabic{enumi})}
\tightlist
\item
  我们取 G′ 为张量积 A′ ⊗\_A G ⊗\_B B′ （第 III 章，第 2 版，附录 II，第 9 小节），并取 u 为映射
\end{enumerate}

\[
g \to 1 \otimes g \otimes 1\tag{\(g\in G\)}
\]

它把 G 映到 G′ 中。如此定义的偶 (G′, u) 显然在下列意义下是泛的：对每个 (A′, B′)-双模 G₁′ 以及每个满足 (11) 的类似条件的、G 到 G₁′ 中的 Z-线性映射 u₁，存在唯一的 G′ 到 G₁′ 中的 Z-线性映射 f，使得 f(a'g'b') = a'f(g')b' （a′ ∈ A′, g′ ∈ G′, b′ ∈ B′；换言之，f 是 G′ 与 G₁′ 的双模结构之间的一个同态），且使得 u₁ = f ∘ u。

\begin{enumerate}
\def\labelenumi{\arabic{enumi})}
\setcounter{enumi}{1}
\item
  当 A = B = G（A 的 (A, A)-双模结构由左、右位似定义）、\(A' = B'\) 且 \(h = h'\) 时，可取 \(G'\) 为环 \(A'\) ，并取 u 为 A 到 \(A'\) 中的同态 h。
\item
  设我们有 A = B、\(A' = B'\) 、\(h = h'\) ，环 A 与 \(A'\) 都是交换的，且 G 的左 A-模结构与其右 A-模结构一致。这时可取 \(G'\) 为张量积 \(A' \otimes_{A} G\) （右 \(A'\) -模结构：即 \(G'\) 的这一结构与其左 \(A'\) -模结构相一致），并取 u 为映射 \(g \to 1 \otimes g\) （ \(g \in G\) ），它把 G 映到 \(G'\) 中。这时我们说，由命题 1（相应地命题 2）所定义的双线性（相应地半双线性）映射 \(\Phi'\) 是由 \(\Phi\) 经基环的扩张、亦即标量扩张而得到的。
\end{enumerate}

以下内容对双线性映射与半双线性映射同样有效；假设与记号即命题 1（相应地命题 2）者。给定 E 或 F 的一个子模 M，我们用 M′ 表示由 M 的典范像所生成的 E′ 或 F′ 的子模。

命题 3. --- 在命题 1（相应地命题 2）的假设与记号下，再设 A、B、A′、B′ 都是体，且映射 \(\alpha\) 与 \(\beta\) ——它们把 A′⊗A G 与 G⊗B B′ 映到 G′ 中，分别由 \(\alpha(a'\otimes g)=a'u(g)\) 与 \(\beta(g\otimes b')=u(g)b'\) （ \(a'\in A'\) ， \(b'\in B'\) ， \(g\in G\) ）刻画——都是单射。设 M 是 E 的一个子空间，N 是 F 的一个子空间。则子空间 \((M')^{0}\) （F′ 中关于 \(\Phi'\) 与 M′ 正交者）等于 \((M^{0})^{\prime}\) ，类似地有 \((N')^{0} = (N^{0})^{\prime}\) 。

事实上，包含关系 \((\mathbf{M}^{0})^{\prime}\subset(\mathbf{M}^{\prime})^{0}\) 与 \((\mathbf{N}^{0})^{\prime}\subset(\mathbf{N}^{\prime})^{0}\) 是显然的（而且对 A、B、A′、B′、\(\alpha\) 或 \(\beta\) 不作任何假设时也为真）。我们将证明包含关系 \((\mathbf{M}^{\prime})^{0}\subset(\mathbf{M}^{0})^{\prime}\) ；我们把完全类似的包含关系 \((\mathbf{N}^{\prime})^{0}\subset(\mathbf{N}^{0})^{\prime}\) 的验证留给读者。设 \(y^{\prime}\) 是 \((\mathbf{M}^{\prime})^{0}\) 的一个元素。可写

\[
y ^ {\prime} = \sum_ {i = 1} ^ {s} y _ {i} \otimes b _ {i} ^ {\prime} \quad (\text { resp. } y ^ {\prime} = \sum_ {i = 1} ^ {s} b _ {i} ^ {\prime} \otimes y _ {i})
\]

其中 \(y_i \in \mathbf{F} (1 \leqslant i \leqslant s)\) ，而诸 \(b_i'\) 是 \(\mathbf{B}'\) 的元素，在 \(\mathbf{B}\) 上线性无关（对于 \(\mathbf{B}\) 在 \(\mathbf{B}'\) 上的左（相应地右）模结构）。设 \(x \in \mathbf{M}\) 与 \(x' = 1 \otimes x \in \mathbf{M}'\) 。我们有

\[
0 = \Phi^ {\prime} (x ^ {\prime}, y ^ {\prime}) = \sum_ {i} u (\Phi (x, y _ {i})) b _ {i} ^ {\prime} = \beta (\sum_ {i} \Phi (x, y _ {i}) \otimes b _ {i} ^ {\prime})
\]

\[
(\text { resp. } 0 = \Phi^ {\prime} (x ^ {\prime}, y ^ {\prime}) = \sum_ {i} u (\Phi (x, y _ {i})) b _ {i} ^ {\prime \mathrm{I}} = \beta (\sum_ {i} \Phi (x, y _ {i}) \otimes b _ {i} ^ {\prime \mathrm{I}}).
\]

由于 β 是单射，且诸 \(b_i'\) （相应地诸 \(b _ {i} ^ {\prime \mathrm{I}}\) ，把 (15) 考虑在内）在 B′ 的左 B-模结构下于 B 上线性无关，这就导出 \(\Phi(x, y_i) = 0\) （对 \(i = 1, \ldots, s\)）。由于这最后一式对每个 \(x \in \mathbf{M}\) 都成立，故有 \(y_i \in \mathbf{M}^0\) （对 \(i = 1, \ldots, s\)），从而 \(y' \in (\mathbf{M}^0)'\) 。证毕。

推论. --- 在命题 3 的假设与记号下，\(\Phi'\) 非退化必须且只需 \(\Phi\) 非退化。

事实上，由命题 3，我们有 \((\mathbf{F}')^0 = (\mathbf{F}^0)'\) 与 \((\mathbf{E}')^0 = (\mathbf{E}^0)'\) 。另一方面，\(\Phi\) （相应地 \(\Phi'\) ）非退化必须且只需 \(\mathbf{F}^0 = \mathbf{E}^0 = \{0\}\) （相应地 \((\mathbf{F}')^0 = (\mathbf{E}')^0 = \{0\}\) ）。

注记. --- 设 A、B、A′ 与 B′ 都是体。那么，对上面三个例子中所定义的双模 G′，映射 α 与 β 都是单射，这由第 III 章，第 2 版，附录 II，第 6 小节立即得出。

\subsection{若干恒等式。}

本节中，我们用 A 表示一个配备了反自同构 J 的环，用 E 表示左 A-模，用 G 表示 (A, A)-双模，并用 Φ 表示 E × E 到 G 中的（关于 J 的）半双线性映射。我们令 Q(x) = Φ(x, x)（x ∈ E）。显然有

\[
\mathrm{Q} (x + y) = \mathrm{Q} (x) + \Phi (x, y) + \Phi (y, x) + \mathrm{Q} (y)\tag{17}
\]

\[
\mathrm{Q} (x - y) = \mathrm{Q} (x) - \Phi (x, y) - \Phi (y, x) + \mathrm{Q} (y)\tag{18}
\]

对 E 中任意的 x, y。于是，相减即得

\[
2 (\Phi (x, y) + \Phi (y, x)) = Q (x + y) - Q (x - y).\tag{19}
\]

设 \(a\) 是 A 的一个元素；在 (19) 中把 \(y\) 换成 \(ay\) ，便得到

\[
2 (\Phi (x, y) a ^ {j} + a \Phi (y, x)) = Q (x + a y) - Q (x - a y).\tag{20}
\]

由 (19) 与 (20)，把 (19) 左乘以 \(a\) 再相减，便得到：

\[
\begin{array}{r l} 2 (a \Phi (x, y) - \Phi (x, y) a ^ {J}) & \\ = a Q (x + y) - a Q (x - y) - Q (x + a y) + Q (x - a y). \end{array}\tag{21}
\]

特别地，设 A 是交换环 K 的一个二次扩张 K(i)，\(i^{2} = -1\) （第 II 章，第 7 节，第 7 小节），J 是 A 的 K-自同构 \(u + iv \rightarrow u - iv\) （u, v 在 K 中），且 G 的左 A-模结构与右 A-模结构相一致。在 (21) 中令 a = i，便得到

\[
4 \Phi (x, y) = Q (x + y) - Q (x - y) + i Q (x + i y) - i Q (x - i y). \tag {22}
\]

\subsection{双线性型与半双线性型。秩。}

本节中，A 表示一个环（相应地一个配备了反自同构 J 的环），E 表示左 A-模，F 表示右 A-模（相应地左 A-模）。我们用左位似与右位似给 A 赋予 (A, A)-双模结构。这时，E × F 到双模 A 中的双线性（相应地关于 J 的右半双线性）映射，称为 E × F 上的双线性（相应地关于 J 的右半双线性）型。

当 E = F 时（这就意味着它是半双线性型），人们常把 E × F 上的一个半双线性型说成 E 上的一个半双线性型。

给定两个左 A-模 E 与 E′ 以及 E 与 E′ 上（关于 J）的两个半双线性型 \(\Phi\) 与 \(\Phi'\) ，称 \(\Phi\) 与 \(\Phi'\) 等价，如果存在 A-模 E 到 A-模 E′ 上的一个同构 \(u\) ，使得 \(\Phi'(u(x), u(y)) = \Phi(x, y)\) 对 E 中一切的 \(x\) 、\(y\) 成立；这时 \(\Phi\) 是 \(\Phi'\) 关于 \(u\) 与 \(u\) 的原像，而 \(\Phi'\) 是 \(\Phi\) 关于 \(u^{-1}\) 与 \(u^{-1}\) 的原像（第 2 小节）。

设 \(\Phi\) 为 \(E \times F\) 上的一个双线性型（其中 F 表示一个右 A-模）。线性映射 \(s_{\Phi}\) 与 \(d_{\Phi}\) （与 \(\Phi\) 相伴者，第 1 小节，定义 2）这时分别是 E 到 F 的对偶 \(F^{*}\) 中与 F 到 E 的对偶 \(E^{*}\) 中的映射。

于是，按定义，

\[
\Phi (x, y) = \left\langle x, d _ {\Phi} (y) \right\rangle = \left\langle y, s _ {\Phi} (x) \right\rangle .\tag{23}
\]

我们现在来定义与半双线性型相伴的线性映射。设 J 是 A 的一个反自同构，\(\Phi\) 是 E × F 上的一个（关于 J 的）（右）半双线性型（其中 F 表示一个左 A-模）；令 J′ = J \(^{-1}\) 。由下式定义的、F × E 到 A 中的映射 \(\Phi'\)

\[
\Phi^ {\prime} (y, x) = \Phi (x, y) ^ {J ^ {\prime}}
\]

\[
(x \in \mathrm{E}, y \in \mathrm{F})
\]

如容易看出的那样，是 F × E 上的一个（关于 J′ 的）（右）半双线性型。由第 2 小节（定义 5），半双线性型 Φ 与 Φ′ 分别等同于 E × F′ 上与 F × E′ 上的双线性型。与这两个双线性型相伴的映射 dΦ 与 dΦ′，分别称为半双线性型 \(\Phi\) 的右相伴与左相伴映射，记作 \(d_{\Phi}\) 与 \(s_{\Phi}\) 。于是，按定义我们有：

\[
\Phi (x, y) = \left\langle x, d _ {\Phi} (y) \right\rangle = \left\langle y, s _ {\Phi} (x) \right\rangle^ {J} \quad (x \in \mathrm{E}, y \in \mathrm{F}).\tag{24}
\]

于是，\(d_{\Phi}\) （相应地 \(s_{\Phi}\) ）是 \(\mathbf{F}^{\mathrm{j}}\) 到 \(\mathbf{E}^{*}\) 中（相应地 \(\mathbf{E}^{\mathrm{j}'}\) 到 \(\mathbf{F}^{*}\) 中）的线性映射，或者等价地说，是 \(\mathbf{F}\) 到 \(\mathbf{E}^{*}\) （相应地 \(\mathbf{E}\) 到 \(\mathbf{F}^{*}\) ）中的、关于 \(\mathbf{J}\) （相应地 \(\mathbf{J}'\) ）的半线性映射，如果我们把 \(\mathbf{J}\) （相应地 \(\mathbf{J}'\) ）看作环 \(\mathbf{A}^{0}\) （\(\mathbf{A}\) 的反环）到 \(\mathbf{A}\) 上的一个同构，并把 \(\mathbf{F}\) （相应地 \(\mathbf{E}\) ）看作一个右 \(\mathbf{A}^{0}\) -模。

公式 (24) 与第 3 小节的定义 6 立即表明，对 F（相应地 E）的每个子模 N（相应地 M），有

\[
\mathbf {N ^ {0}} = \stackrel {- 1} {s _ {\Phi}} (\mathbf {N ^ {\prime}}) \quad (\text { resp. } \mathbf {M ^ {0}} = \stackrel {- 1} {d _ {\Phi}} (\mathbf {M ^ {\prime}}))\tag{25}
\]

其中 N′（相应地 M′）是 F 的对偶 F*（相应地 E 的对偶 E*）中与 N（相应地 M）正交的子模（第 II 章，第 4 节，第 2 小节）。

命题 4. --- 设 A 是一个体，\(\Phi\) 是 \(E \times F\) 上的一个双线性型（相应地关于 J 的半双线性型）；\(E/F^{0}\) 是有限维的必须且只需 \(F/E^{0}\) 是有限维的，且这时两者的维数相等。

事实上，设 \(\Phi_{1}\) 为与 \(\Phi\) 相伴的非退化型，定义在 \((\mathrm{E} / \mathrm{F}^{0}) \times (\mathrm{F} / \mathrm{E}^{0})\) 上（第 3 小节）。设 \(\mathrm{E} / \mathrm{F}^{0}\) 是有限维 \(n\) ；由于线性映射 \(d_{\Phi_1}\) （从 \(\mathrm{F} / \mathrm{E}^{0}\) （相应地 \((\mathrm{F} / \mathrm{E}^{0})^{\mathrm{j}})\) ）到 \((\mathrm{E} / \mathrm{F}^{0})^{*}\) 中）是单射，\(\mathrm{F} / \mathrm{E}^{0}\) 是有限维 \(n' \leqslant n\) ；对 \(s_{\Phi_1}\) 作同样考虑，同样看出 \(n \leqslant n'\) 。

推论 1. --- 设 A 是一个体且 \(\Phi\) 非退化。E 的子空间 M 是有限维的，必须且只需 M \(^{0}\) 在 F 中具有有限余维数，且这时有 codim M \(^{0}\) = dim M，以及 M \(^{00}\) = M。

由于 \(F^{0}=\{0\}\) ，前两个断言由把命题 4 应用于 \(\Phi\) 在 M × F 上的限制得出。此外，\(M^{0}\) 是 \(M^{00}\) 的正交补，故 \(M^{00}\) 的维数有限且等于 codim \(M^{0}=dim M\) ；但由于 \(M^{00}\supset M\) ，就有 \(M^{00}=M\) 。

推论 2. --- 在推论 1 的假设下，设 M、N 是 E 的两个子空间；则有 \((M + N)^{0} = M^{0} \cap N^{0}\) ；若此外 M 与 N 都是有限维的，则有 \((M \cap N)^{0} = M^{0} + N^{0}\) 。

第一个断言是显然的。设 M 与 N 都是有限维的，令 \(G = M^{0} + N^{0}\) ；由推论 1 有 \(G^{0} = M^{00} \cap N^{00} = M \cap N\) ；把命题 4 应用于 \(\Phi\) 在 \(M \times G\) 上的限制，便表明（由于 \(M^{0} \subset G\) 且 \(G^{0} \subset M\) ）有 dim \(M/(M \cap N) = \dim G/M^{0} = \text{codim } M^{0} - \text{codim } G\) ，而由 codim \(M^{0} = \dim M\) ，便导出 dim \((M \cap N) = \text{codim } G\) 。但由推论 1，同样有 dim \((M \cap N) = \text{codim } (M \cap N)^{0}\) ，而由于 \(G \subset G^{00} = (M \cap N)^{0}\) ，就有 \(G = (M \cap N)^{0}\) 。

命题 4 使我们可以给出下面的定义：

定义 7. --- 设 A 是一个体（相应地一个配备了反自同构 J 的体），E 是 A 上的一个左向量空间，F 是 A 上的一个右（相应地左）向量空间，Φ 是 E × F 上的一个双线性（相应地关于 J 的半双线性）型。设 E/F⁰ 与 F/E⁰ 在 A 上都是有限维的。Φ 的秩就是向量空间 E/F⁰ 与 F/E⁰ 的公共（有限）维数。

当 \(E/F^{0}\) 与 \(F/E^{0}\) 是无限维时，称 \(\Phi\) 具有无穷秩。

命题 5. --- 在定义 7 的假设与记号下，线性映射 \(s_{\Phi}\) 与 \(d_{\Phi}\) （与 \(\Phi\) 相伴者）有相同的秩，且这个秩等于型 \(\Phi\) 的秩。

事实上，F 到 E* 中的映射 \(d_{\Phi}\) 的核显然是 E°，故其秩等于 F/E° 的维数。同样，\(s_{\Phi}\) 的秩等于 E/F° 的维数。

命题 6. --- 在定义 7 的假设与记号下，再设 E 与 F 有相同的有限维数。则下列条件等价：

\begin{enumerate}
\def\labelenumi{\alph{enumi})}
\item
  \(d_{\Phi}\) 是单射；
\item
  \(d_{\Phi}\) 是满射；
\item
  \(s_{\Phi}\) 是单射；
\end{enumerate}

\(d)\) \(s_{\Phi}\) 是满射；

\begin{enumerate}
\def\labelenumi{\alph{enumi})}
\setcounter{enumi}{4}
\tightlist
\item
  \(\Phi\) 非退化。
\end{enumerate}

事实上，由于 E、F、E* 与 F* 有相同的有限维数，a) 与 b) 等价，c) 与 d) 也等价（第 II 章，第 3 节，第 4 小节）。由于 \(s_{\Phi}\) 与 \(d_{\Phi}\) 有相同的秩（命题 5），a) 与 c) 等价。由于 e) 等价于关系式 \(\mathrm{E}^{0} = \mathrm{F}^{0} = \{0\}\) ，它等价于 a) 与 c) 的合取，从而所述诸条件等价。

推论. --- 在定义 7 的假设与记号下，再设 E 是有限维的且 \(\Phi\) 非退化。则有 dim E = dim F，且对 E 的每个基 \((e_i)\) （ \(1 \leqslant i \leqslant \dim E\) ），存在 F 的一个基 \((f_i)\) ，使得 \(\Phi(e_i, f_k) = \delta_{ik}\) （ \(i, k = 1, \ldots, \dim E\) ）。

事实上，由于 \(\Phi\) 非退化，我们有 \(E^{0} = F^{0} = \{0\}\) ，从而 dim \(E = \dim F\) （命题 4）。由此（命题 6），\(d_{\Phi}\) 是 \(F\) （相应地 \(F^{j}\) ）到 \(E^{*}\) 上的一个同构；因此，若 \((e_{i}^{*})\) 是 \((e_{i})\) 的对偶基，则元素 \(f_{i} = d_{\Phi}^{-1}(e_{i}^{*})\) 构成 \(F\) 的一个基，它鉴于公式 (23)（相应地公式 (24)），满足 \(\Phi(e_{i}, f_{k}) = \delta_{ik}\) 。

显然，在这个推论中，可以把 E 与 F 的作用互换，在证明中把 \(d_{\Phi}\) 换成 \(s_{\Phi}\) 。

注记. --- 设 A 是一个配备了反自同构 J 的环，M 与 N 是右 A-模，\(\Phi\) 是 M × N 上的一个（关于 J 的）左半双线性型（第 2 小节，注记）；于是它满足等式

\[
\Phi (x a, y a ^ {\prime}) = a ^ {j} \Phi (x, y) a ^ {\prime} \quad (a, a ^ {\prime} \in A, x \in M, y \in N).
\]

映射 \(\Phi'\) （由 \(\Phi'(y, x) = \Phi(x, y)^{J'}\) 定义，其中 J′ = J \(^{-1}\) ，从 N × M 到 A 中）是一个关于 J′ 的左半双线性型，且 \(\Phi\) 与 \(\Phi'\) 分别等同于 M \(^{J'}\) × N 与 N \(^{J'}\) × M 上的双线性型。与这两个双线性型相伴的映射 \(s_{\Phi}\) 与 \(s_{\Phi'}\) ，分别称为半双线性型 \(\Phi\) 的左相伴与右相伴映射，记作 \(s_{\Phi}\) 与 \(d_{\Phi}\) 。于是，按定义我们有

\[
\Phi (x, y) = \left\langle y, s _ {\Phi} (x) \right\rangle = \left\langle x, d _ {\Phi} (y) \right\rangle^ {J} \quad (x \in M, y \in N),\tag{26}
\]

且 \(s_{\Phi}\) （相应地 \(d_{\Phi}\) ）是 \(\mathbf{M}^{\mathbf{j}}\) 到 \(\mathbf{N}^{*}\) 中的一个线性映射（相应地，

N \(^{j'}\) 到 M* 中的线性映射）。对这里所考虑的情形，容易叙述并证明定义 7 与命题 4、5、6 的类似结果。

\subsection{双线性型或半双线性型的逆型。}

设 A 是一个环，E 是一个左 A-模，F 是一个右 A-模，\(\Phi\) 是 \(\mathbf{E} \times \mathbf{F}\) 上的一个双线性型。这里假设 \(\Phi\) 的相伴映射（将记作 \(s\) 与 \(d\) ）是双射。于是，积映射 \((s, d)\) 是 \(\mathbf{E} \times \mathbf{F}\) 到 \(\mathbf{F}^* \times \mathbf{E}^*\) 上的一个双射，且经由结构转移，定义了双线性型 \(\hat{\Phi}\) （在 \(\mathbf{F}^* \times \mathbf{E}^*\) 上）。因此它满足

\[
\begin{array}{r l} \hat {\Phi} (y ^ {\prime}, x ^ {\prime}) = & \Phi (s ^ {- 1} (y ^ {\prime}), d ^ {- 1} (x _ {i} ^ {\prime})) \\ = & \left\langle s ^ {- 1} (y ^ {\prime}), x ^ {\prime} \right\rangle = \left\langle d ^ {- 1} (x ^ {\prime}), y ^ {\prime} \right\rangle \qquad (x ^ {\prime} \in \mathrm{E} ^ {*}, y ^ {\prime} \in \mathrm{F} ^ {*}). \end{array}\tag{27}
\]

定义 8. --- 设 \(\Phi\) 是 E × F 上的一个双线性型，它的相伴映射 s 与 d 是双射。由 (27) 定义的 F* × E* 上的双线性型 \(\hat{\Phi}\) 称为 \(\Phi\) 的逆型。

现在设 \(\hat{s}\) 与 \(\hat{d}\) 为左（相应地右）相伴于 \(\hat{\Phi}\) 的、F* 到 E** 与 E* 到 F** 中的线性映射。由于对 \(x' \in \mathrm{E}^*\) 与 \(y' \in \mathrm{F}^*\) ，按定义有

\[
\hat {\Phi} (y ^ {\prime}, x ^ {\prime}) = \left\langle y ^ {\prime}, \hat {d} (x ^ {\prime}) \right\rangle = \left\langle x ^ {\prime}, \hat {s} (y ^ {\prime}) \right\rangle
\]

与 (27) 相比较，便看出 F* 上的线性型 \(\hat{d}(x')\) 等于由 F 的元素 \(d^{-1}(x')\) 所定义的线性型。由此得出，复合映射 \(\hat{d} \circ d\) 是 F 到其双对偶 F** 中的典范映射，且这个映射是双射，因为 \(d\) 与 \(\hat{d}\) （后者经由结构转移）都是双射；于是，若把 F 与 F** 典范地等同起来，就有 \(\hat{d} = d^{-1}\) 。类似地，E 与 E** 典范地等同，而 E 到 E** 的典范映射是 \(\hat{s} \circ s\) ，且有 \(\hat{s} = s^{-1}\) 。由此推出 \(\hat{\Phi}\) 的逆型是 \(\Phi\) 。

现在考虑一个配备了反自同构 J 的环 A、两个左 A-模 E 与 F，以及一个（关于 J 的）右半双线性型 \(\Phi\) （在 E \(\times\) F 上），使得 \(\Phi\) 的相伴映射（将记作 \(s\) 与 \(d\) ）是双射。用 (27) 的第一个等式定义一个映射 \(\hat{\Phi}\) （从 \(\mathbf{F}^{*} \times \mathbf{E}^{*}\) 到 A 中）。由 (24)（第 6 小节），这个映射满足关系式

\[
\hat {\Phi} (y ^ {\prime}, x ^ {\prime}) = \left\langle s ^ {- 1} (y ^ {\prime}), x ^ {\prime} \right\rangle = \left\langle d ^ {- 1} (x ^ {\prime}), y ^ {\prime} \right\rangle^ {J} \quad (x ^ {\prime} \in E ^ {*}, y ^ {\prime} \in F ^ {*}).\tag{28}
\]

映射 \(\hat{\Phi}\) 显然是 \(\mathbf{Z}\) -双线性的；此外，对 A 中的 \(a\) 、\(b\) 、\(x' \in \mathrm{E}^*\) 与 \(y' \in \mathrm{F}^*\) ，鉴于 \(s\) 与 \(d\) 的定义，我们有

\[
\widehat {\Phi} (y ^ {\prime} a, x ^ {\prime} b) = \Phi (a ^ {j} s ^ {- 1} (y ^ {\prime}), b ^ {j ^ {\prime}} d ^ {- 1} (x ^ {\prime})) = a ^ {j} \widehat {\Phi} (y ^ {\prime}, x ^ {\prime}) b;
\]

因此 \(\hat{\Phi}\) 是 \(\mathbf{F}^{*} \times \mathbf{E}^{*}\) 上的一个（关于 J 的）左半双线性型（第 2 小节）。

定义 9. --- 设 \(\Phi\) 是 E × F 上的一个（关于 J 的）右半双线性型，其相伴映射 s 与 d 是双射。F* × E* 上的（关于 J 的）左半双线性型 \(\hat{\Phi}\) 称为 \(\Phi\) 的逆型。

我们把左半双线性型的逆型的定义与研究留给读者。这个逆型是一个右半双线性型。

设 \(\hat{s}\) 与 \(\hat{d}\) 为 \(\hat{\Phi}\) 的相伴映射；由 (26)（第 6 小节），我们有

\[
\hat {\Phi} (y ^ {\prime}, x ^ {\prime}) = \left\langle y ^ {\prime}, \hat {d} (x ^ {\prime}) \right\rangle^ {\mathbf {J}} = \left\langle x ^ {\prime}, \hat {s} (y ^ {\prime}) \right\rangle .\tag{29}
\]

由于 \(s\) 是双射，并由于由 (28) 与 (29) 得出的等式 \(\langle s^{-1}(y'), x'\rangle = \langle y', \hat{d}(x') \rangle^{\mathfrak{J}}\) ，我们导出 \(\hat{d}\) 是双射；于是 \(\hat{d} \circ d\) 是双射。但由 (28) 与 (29) 得出的等式 \(\langle d^{-1}(x'), y' \rangle = \langle y', \hat{d}(x') \rangle\) 表明 \(\hat{d} \circ d\) 是 F 到其双对偶 F** 中的典范映射。同样看出 \(\hat{s}\) 是双射且 \(\hat{s} \circ s\) 是 E 到 E** 中的典范映射。因此，若借助这些典范映射把 E** 与 E、F** 与 F 等同起来，我们就有 \(\hat{s} = s^{-1}\) 、\(\hat{d} = d^{-1}\) ，且 \(\Phi\) 是 \(\hat{\Phi}\) 的逆型。

在同一记号与假设下，设 a 是 A 的中心的一个可逆元素。于是，由 (23)（相应地 (24)），与型 \(a\Phi\) 相伴的映射等于 \(a.d\) 与 \(a.s\) （相应地 \(a^{j'}\) .s），从而是双射。于是由 (27) 推出，\(a\Phi\) 的逆型是 \(a^{-1}\widehat{\Phi}\) （相应地 \((a^{j'})^{-1}\widehat{\Phi}\) ）。

\subsection{同态的伴随。}

本节中，我们用 A 表示一个环（相应地一个配备了反自同构 J 的环），用 E 与 E′ 表示两个左 A-模，用 F 与 F′ 表示两个右 A-模（相应地左 A-模），并用 Φ 与 Φ′ 分别表示 E × F 与 E′ × F′ 上的两个双线性型（相应地关于 J 的半双线性型）。我们假设 Φ 非退化，换言之（第 1 小节），与 Φ 相伴的线性映射 \(d_{\Phi}\) 与 \(s_{\Phi}\) 都是单射。

给定一个同态 \(u\) （从 E 到 \(\mathbf{E}'\) 中），考虑这样的集合 \(\mathbf{F}_{1}'\) ：它的元素 \(y'\) 属于 \(\mathbf{F}'\) ，且存在 \(y \in \mathbf{F}\) 使得有 \(d_{\Phi'}(y') \circ u = d_{\Phi}(y)\) ，也就是说 \(\Phi'(u(x), y') = \Phi(x, y)\) （对一切 \(x \in \mathbf{E}\) ）。显然 \(\mathbf{F}_{1}'\) 是 \(\mathbf{F}'\) 的一个子模。由于 \(d_{\Phi}\) 是单射，对每个 \(y' \in \mathbf{F}_{1}'\) ，存在且仅存在 F 的一个元素 \(y\) ，使得 \(\Phi'(u(x), y') = \Phi(x, y)\) 。如此定义的映射 \(y' \to y\) （从 \(\mathbf{F}_{1}'\) 到 F 中）是 A-线性的；把它记作 \(u^{*}\) ，则对每个 \(x \in \mathbf{E}\) 与每个 \(y' \in \mathbf{F}_{1}'\) ，有

\[
\Phi^ {\prime} (u (x), y ^ {\prime}) = \Phi (x, u ^ {*} (y ^ {\prime})).\tag{30}
\]

定义 10. --- 在上述假设与记号下，称同态 \(u^*\) （从 \(F_1'\) 到 F 中，满足 (30)）为 \(u\) 的左伴随，而称 \(F_1'\) 为 \(u^*\) 的定义子模。

类似地，用公式定义 F 到 F′ 中的同态 \(\nu\) 的右伴随

\[
\Phi^ {\prime} (x ^ {\prime}, \nu (y)) = \Phi (\nu^ {*} (x ^ {\prime}), y) \quad (x ^ {\prime} \in \mathrm{E} _ {1} ^ {\prime}, y \in \mathrm{F}),\tag{31}
\]

其中 \(E_{1}^{\prime}\) 表示 \(E^{\prime}\) 的、与 \(F_{1}^{\prime}\) 类似地定义的子模。

注记. --- 如果 \(u^{*}\) （\(u: E \to E'\) 的左伴随）处处有定义，且 \(s_{\Phi'}\) 与 \(d_{\Phi'}\) 都是单射，则公式 (30) 表明 \(u\) 是 \(u^{*}\) 的右伴随。

由 (30) 我们导出：若 \(u_1\) 与 \(u_2\) 是 E 到 E′ 中的两个同态，其伴随都处处有定义，且 c 是 A 的中心的一个元素，则有

\[
\left\{ \begin{array}{l} (u _ {1} + u _ {2}) ^ {*} = u _ {1} ^ {*} + u _ {2} ^ {*}; 1 ^ {*} = 1; \\ (c u _ {1}) ^ {*} = c. u _ {1} ^ {*} \quad \text { when } \Phi \text { and } \Phi^ {\prime} \text { are bilinear }; \\ (c u _ {1}) ^ {*} = c ^ {J ^ {\prime}}. u _ {1} ^ {*} \quad \text { when } \Phi \text { and } \Phi^ {\prime} \text { are sesquilinear }. \end{array} \right.\tag{32}
\]

此外，若 \(E^{\prime\prime}\) 是第三个左 A-模，\(F^{\prime\prime}\) 是第三个右（相应地左）A-模，\(\Phi^{\prime\prime}\) 是 \(E^{\prime\prime} \times F^{\prime\prime}\) 上的一个双线性型（相应地关于 J 的半双线性型），且 \(u^{\prime}\) 是 \(E^{\prime}\) 到 \(E^{\prime\prime}\) 中的一个同态，其（左）伴随处处有定义，则有

\[
(u ^ {\prime} \circ u) ^ {*} = u ^ {*} \circ u ^ {\prime *}.\tag{33}
\]

特别地，若 \(u\) 是 E 到 E′ 上的一个同构，且伴随 \(u^{*}\) 与 \((u^{-1})^{*}\) 都处处有定义，则 \(u^{*}\) 是 F′ 到 F 上的一个同构，且我们有 \((u^{*})^{-1} = (u^{-1})^{*}\) 。右伴随有类似的性质。

命题 7. --- 在与上面相同的记号下，设 \(d_{\Phi}\) 是双射。于是每个 E 到 E′ 中的同态 \(u\) 都有一个处处有定义的左伴随，且我们有 \(u^{*} = (d_{\Phi})^{-1} \circ {}^{t}u \circ d_{\Phi'}\) 。

事实上，由于 \(d_{\Phi}\) 是双射，用本节开头的记号，有 \(\mathbf{F}_{1}^{\prime}=\mathbf{F}^{\prime}\) ，因而 \(u^{*}\) 处处有定义。另一方面，(30) 等价于

\[
\left\langle u (x), d _ {\Phi^ {\prime}} \left(y ^ {\prime}\right) \right\rangle = \left\langle x, \left(d _ {\Phi} \circ u ^ {*}\right) \left(y ^ {\prime}\right) \right\rangle \quad \left(x \in \mathrm{E}, y ^ {\prime} \in \mathrm{F} ^ {\prime}\right);
\]

或者 \(\langle d_{\Phi'}(y'), u(x) \rangle = \langle {}^t u(d_{\Phi'}(y')), x \rangle\) ；于是 \({}^t u(d_{\Phi'}(y')) = d_{\Phi}(u^*(y'))\) （对一切 \(y' \in F'\) ），从而 \({}^t u \circ d_{\Phi'} = d_{\Phi} \circ u^*\) ，因此得到所宣布的 \(u^*\) 的表达式。证毕。

注记. --- 当 \(s_{\Phi}\) 是双射时，每个 F 到 F′ 中的同态 \(\nu\) 都有一个处处有定义的右伴随，且我们有

\[
v ^ {*} = (s _ {\Phi}) ^ {- 1} \circ {} ^ {t} v \circ s _ {\Phi^ {\prime}}.\tag{34}
\]

命题 8. --- 在与上面相同的记号下，设 \(s_{\Phi}\) 与 \(d_{\Phi}\) 都是双射。设 \(u\) 与 \(v\) 分别是 E 到 \(\mathbf{E}'\) 上与 F 到 \(\mathbf{F}'\) 上的同构。于是，\(\Phi\) 是 \(\Phi'\) 关于 \(u\) 与 \(\nu\) 的原像（也就是说，有 \(\Phi(x,y)=\Phi'(u(x),\nu(y))\) （对一切 \(x\in\mathbf{E},y\in\mathbf{F}\) ））必须且只需我们有 \(u^{-1}=\nu^{*}\) 与 \(\nu^{-1}=u^{*}\) 。

事实上，\(\Phi'(u(x), v(y)) = \Phi(x, y)\) 也可写成 \(\Phi(x, u^*(v(y))) = \Phi(x, y)\) 。若它对一切 \(x \in E\) 与 \(y \in F\) 成立，则有 \(u^* \circ v = 1\) （由于 \(\Phi\) 非退化）。于是由 (33)，同样有 \(v^* \circ u = 1\) 。其逆是显然的。

推论. --- 设 A 是一个配备了反自同构 J 的环，E 是一个左 A-模，\(\Phi\) 是 E × E 上的一个（关于 J 的）半双线性型，其相伴映射都是双射，u 是 A-模 E 的一个自同构。u 保持 \(\Phi\) 不变（也就是说，对 E 中一切的 x, y 有 \(\Phi(u(x), u(y)) = \Phi(x, y)\) ）必须且只需 u 的两个伴随相等，且我们有 \(u^{*} = u^{-1}\) 。

这由命题 8 立即得出。

注记. --- 在命题 8 的推论的假设下，再设 A 是一个体且 E 在 A 上是有限维的。设 \(w\) 是 E 的一个自同态，\(w_{1}\) 与 \(w_{2}\) 是它的右伴随与左伴随。条件 \(ww_{1}=1\) 、\(ww_{2}=1\) 、\(w_{1}w=1\) 、\(w_{2}w=1\) 中的每一个都蕴涵 \(w\) 是 E 的一个保持 \(\Phi\) 不变的自同构，且 \(w_{1}=w_{2}\) 。

\subsection{半双线性型的张量积与外幂。}

本节中，A 表示一个交换环。因此，两个 A-模的积上的双线性型是半双线性型的一个特例。我们用 J 表示 A 的一个自同构，用 J′ 表示其逆。

设 \(\mathbf{E}_i\) （ \(i = 1, \ldots, m\) ）是一些 A-模。映射

\[
(x _ {1}, \dots , x _ {m}) \rightarrow x _ {1} \otimes \dots \otimes x _ {m}
\]

（从 \(\prod_{i=1}^{m}\mathrm{E}_{i}^{j}\) 到 \((\bigotimes_{i=1}^{m}\mathrm{E}_{i})^{j}\) 中，\(x_{i}\in\mathrm{E}_{i}^{j}\) ）（参看第 2 小节，定义 5）显然是 A-重线性的；于是它定义了（第 III 章，第 1 节，第 7 小节）一个 A-线性映射 \(f\) （从 \(\bigotimes_{i}\mathrm{E}_{i}^{\mathrm{J}}\) 到 \((\bigotimes_{i}\mathrm{E}_{i})^{\mathrm{J}}\) 中）；这个映射把 \(x_{1}\otimes\cdots\otimes x_{m}\) （其中记号 \(\otimes\) 表示 \(\bigotimes_{i}\mathrm{E}_{i}^{\mathrm{J}}\) 中的张量积）映为 \(x_{1}\otimes\cdots\otimes x_{m}\) （其中记号 \(\otimes\) 表示 \((\bigotimes_{i}\mathrm{E}_{i})^{\mathrm{J}}\) 中的张量积）。于是 \(f\) 是 \(\bigotimes_{i}\mathrm{E}_{i}^{\mathrm{J}}\) 到 \((\bigotimes_{i}\mathrm{E}_{i})^{\mathrm{J}}\) 上的一个同构。我们将借助这个同构把这两个模等同起来。

类似地，设 E 是一个 A-模。映射

\[
(x _ {1}, \dots , x _ {m}) \rightarrow x _ {1} \wedge \dots \wedge x _ {m}
\]

（从 \((\mathrm{E}^{\mathfrak{J}})^{m}\) 到 \((\wedge^{\mathfrak{m}}\mathrm{E})^{\mathfrak{J}}\) 中）显然是 A-重线性的且交错的。于是它定义了一个 A-线性映射 \(f\) （从 \(\wedge^{\mathfrak{m}}\mathrm{E}^{\mathfrak{J}}\) 到 \((\wedge^{\mathfrak{m}}\mathrm{E})^{\mathfrak{J}}\) 中），后者显然是一个同构。我们将借助这个同构把 \(\wedge^{\mathfrak{m}}\mathrm{E}^{\mathfrak{J}}\) 与 \((\wedge^{\mathfrak{m}}\mathrm{E})^{\mathfrak{J}}\) 等同起来。

设 \(x'\) 是 \(\mathbf{E}^*\) （\(\mathbf{E}\) 的对偶）的一个元素。映射 \(x \to \langle x, x' \rangle^{\mathfrak{J}}\) （ \(x \in \mathbf{E}\) ）是一个元素 \(x'^{\mathfrak{J}}\) （属于 \((\mathrm{E}^{\mathfrak{J}})^*\) ），且立即看出 \(x' \to x'^{\mathfrak{J}}\) 是一个双射 \(g\) （从 \(\mathbf{E}^*\) 到 \((\mathrm{E}^{\mathfrak{J}})^*\) 上），它满足 \(g(ax') = a^{\mathfrak{J}}g(x')\) （对一切 \(a \in \mathbf{A}\) ）。因此，\(g\) 与 \((\mathrm{E}^*)^{\mathfrak{J}}\) 到 \(\mathbf{E}^*\) 上的恒同映射的复合，是 \((\mathrm{E}^*)^{\mathfrak{J}}\) 到 \((\mathrm{E}^{\mathfrak{J}})^*\) 上的一个同构。我们将借助这个同构把这些模等同起来，并记作 \(\mathbf{E}_j^*\) 。

设 \(E_{i}, F_{i} (i = 1, \ldots, m)\) 是一些 A-模，\(\Phi_{i} (i = 1, \ldots, m)\) 是 \(E_{i} \times F_{i}\) 上的（关于 J 的）半双线性型。映射

\[
(x _ {1}, \dots , x _ {m}, y _ {1}, \dots , y _ {m}) \rightarrow \Phi_ {1} (x _ {1}, y _ {1}) \Phi_ {2} (x _ {2}, y _ {2}) \dots \Phi_ {m} (x _ {m}, y _ {m})
\]

\((x_{i} \in \mathrm{E}_{i}, y_{i} \in \mathrm{F}_{i}, i = 1, \ldots, m)\) 是 \(\mathrm{E}_{1} \times \cdots \times \mathrm{E}_{m} \times \mathrm{F}_{1}^{J} \times \cdots \times \mathrm{F}_{m}^{J}\) 到 A 中的一个 A-重线性映射，从而定义了 \((\bigotimes_{i} \mathrm{E}_{i}) \times (\bigotimes_{i} \mathrm{F}_{i}^{J})\) 上的一个双线性型（第 III 章，第 1 节，第 7 小节）。由于第二个因子已与 \((\bigotimes_{i} \mathrm{F}_{i})^{J}\) 等同，我们便定义了一个（关于 J 的）半双线性型 \(\Phi\) （在 \((\bigotimes_{i} \mathrm{E}_{i}) \times (\bigotimes_{i} \mathrm{F}_{i})\) 上）。它由下式刻画：

\[
\Phi (x _ {1} \otimes \dots \otimes x _ {m}, y _ {1} \otimes \dots \otimes y _ {m}) = \prod_ {i = 1} ^ {m} \Phi_ {i} (x _ {i}, y _ {i}) \quad (x _ {i} \in \mathrm{E} _ {i}, y _ {i} \in \mathrm{F} _ {i}).\tag{35}
\]

定义 11. --- 给定 A-模 \(\mathbf{E}_{i}, \mathbf{F}_{i}\) （ \(i=1,\ldots,m\) ），并对每个 \(i\) 给定一个（关于 J 的）半双线性型 \(\Phi_{i}\) （在 \(\mathbf{E}_{i} \times \mathbf{F}_{i}\) 上），由 (35) 刻画的半双线性型 \(\Phi\) （关于 J，在 \((\bigotimes_{i}\mathbf{E}_{i}) \times (\bigotimes_{i}\mathbf{F}_{i})\) 上），称为诸半双线性型 \(\Phi_{i}\) 的张量积。

在诸 \(E_{i}\) 与诸 \(F_{i}\) 都等于同一个模 E 且诸 \(\Phi_{i}\) 都等于同一个型 \(\Psi\) 的情形，称 \(\Phi\) 为 \(\Psi\) 到 \(\otimes^{m}E\) 上的扩张。

在定义 11 的记号下，我们来研究与 \(\Phi\) 相伴的映射。由公式 (24)（第 6 小节）与 (35)，我们导出关系式

\[
\Phi (x _ {1} \otimes \dots \otimes x _ {m}, y _ {1} \otimes \dots \otimes y _ {m}) = \prod_ {i = 1} ^ {m} \left\langle x _ {i}, d _ {\Phi_ {i}} (y _ {i}) \right\rangle = \prod_ {i = 1} ^ {m} \left\langle y _ {i}, s _ {\Phi_ {i}} (x _ {i}) \right\rangle^ {j}.
\]

因此我们有：

\[
s _ {\Phi} = j _ {s} \circ (s _ {\Phi_ {1}} \otimes \dots \otimes s _ {\Phi_ {m}}), \quad d _ {\Phi} = j _ {d} \circ (d _ {\Phi_ {1}} \otimes \dots \otimes d _ {\Phi_ {m}})\tag{36}
\]

其中 \(j_s\) （相应地 \(j_d\) ）表示 \(\bigotimes_i F_i^*\) 到 \((\bigotimes_i F_i)^*\) 中（相应地 \(\bigotimes_i E_i^*\) 到 \((\bigotimes_i E_i)^*\) 中）的典范映射（第 III 章，第 1 节，第 4 与第 7 小节）。

命题 9. — 设 A 是一个配备了自同构 J 的交换域，E \(_{i}\) 、F \(_{i}\) 是 A 上两个有限维向量空间，Φ \(_{i}\) 是 E \(_{i}\) × F \(_{i}\) 上的一个（关于 J 的）半双线性型（1 ≤ i ≤ m）。如果诸型 Φ \(_{i}\) 都非退化，则它们的张量积 Φ 也非退化。这时 \(\hat{\Phi}\) （\(\Phi\) 的逆型）是诸逆型 Φ \(_{i}\) 的张量积。

事实上，由于 A 是一个体，由第 III 章，第 1 节，第 3 小节的命题 6 与命题 7 可知，A-模的单射（相应地满射）线性映射的张量积是单射（相应地满射）线性映射。由于诸 \(s_{\Phi_i}\) 依假设是双射（命题 6，第 6 小节），故它们的张量积也是。另一方面，典范映射 \(j_s\) （从 \(\bigotimes_i F_i^*\) 到 \((\bigotimes_i F_i)^*\) 中）是双射（第 III 章，第 1 节，第 5 小节，命题 11）。因此，鉴于 (36)，\(s_\Phi\) 是双射，这就建立了我们的第一个断言（命题 6，第 6 小节）。同样，\(d_\Phi\) 也是双射。

在第二个断言中，我们已把 \(\bigotimes_{i} \mathrm{F}_{i}^{*}\) 与 \((\bigotimes_{i} \mathrm{F}_{i})^{*}\) 、\(\bigotimes_{i} \mathrm{E}_{i}^{*}\) 与 \((\bigotimes_{i} \mathrm{E}_{i})^{*}\) 借助映射 \(j_{s}\) 与 \(j_{d}\) （它们在此是同构）隐式地等同起来。陈述中所引用的诸逆型存在，因为诸 \(s_{\Phi_{i}}\) 、诸 \(d_{\Phi_{i}}\) 、\(s_{\Phi}\) 与 \(d_{\Phi}\) 都是双射（第 7 小节）。令 \(x' = x_{1}' \otimes \cdots \otimes x_{m}'\) ，\(y' = y_{1}' \otimes \cdots \otimes y_{m}'\) （ \(x_{i}' \in \mathrm{E}_{i}^{*}\) ，\(y_{i}' \in \mathrm{F}_{i}^{*}\) ，\(i = 1, \ldots, m\) ）。由逆型的定义，并鉴于 (36)，我们有

\[
\begin{array}{l} \widehat {\Phi} (j _ {s} (y ^ {\prime}), j _ {d} (x ^ {\prime})) = \Phi (s _ {\Phi_ {1}} ^ {- 1} (y _ {1} ^ {\prime}) \otimes \dots \otimes s _ {\Phi_ {m}} ^ {- 1} (y _ {m} ^ {\prime}), d _ {\Phi_ {1}} ^ {- 1} (x _ {1} ^ {\prime}) \otimes \dots \otimes d _ {\Phi_ {m}} ^ {- 1} (x _ {m} ^ {\prime})) \\ = \prod_ {i = 1} ^ {m} \Phi_ {i} (s _ {\Phi_ {i}} ^ {- 1} (y _ {i} ^ {\prime}), d _ {\Phi_ {i}} ^ {- 1} (x _ {i} ^ {\prime})) = \prod_ {i = 1} ^ {m} \widehat {\Phi} _ {i} (y _ {i} ^ {\prime}, x _ {i} ^ {\prime}), \end{array}
\]

由此即得我们的第二个断言。

证毕。

设 E 与 F 是交换环 A 上两个模，\(\Phi\) 是 E × F 上的一个（关于 J 的）半双线性型。映射

\[
(x _ {1}, \dots , x _ {m}, y _ {1}, \dots , y _ {m}) \rightarrow \det (\Phi (x _ {i}, y _ {k})) \quad (x _ {i} \in \mathrm{E}, y _ {i} \in \mathrm{F}, i = 1, \dots , m)
\]

是 \(\mathbf{E}^{m}\times(\mathbf{F}^{\mathrm{j}})^{m}\) 到 A 中的 A-重线性映射。因此它定义了一个双线性型 \(\Phi'\) （在 \((\bigotimes^{\overline{m}}\mathbf{E})\times(\bigotimes^{\overline{m}}\mathbf{F}^{\mathrm{j}})\) 上），由下式刻画

\[
\Phi^ {\prime} (x _ {1} \otimes \dots \otimes x _ {m}, y _ {1} \otimes \dots \otimes y _ {m}) = \det (\Phi (x _ {i}, y _ {k})).
\]

由于当 \(x_{i}=x_{k}\) 或 \(y_{i}=y_{k}\) （ \(i\neq k\) ）时第一项为零，\(\Phi^{\prime}\) 通过过渡到商，定义了 \((\bigwedge^{m}\mathrm{E})\times(\bigwedge^{m}\mathrm{F}^{\mathrm{J}})\) 上的一个双线性型，或者再说一遍，由于 \(\bigwedge^{m}\mathrm{F}^{\mathrm{J}}\) 与 \((\bigwedge^{m}\mathrm{F})^{\mathrm{J}}\) 等同，它定义了一个（关于 J 的）半双线性型 \(\Phi_{(m)}\) （在 \((\bigwedge^{m}\mathrm{E})\times(\bigwedge^{m}\mathrm{F})\) 上）。它由下式刻画

\[
\left\{ \begin{array}{l} \Phi_ {(m)} (x _ {1} \wedge \dots \wedge x _ {m}, y _ {1} \wedge \dots \wedge y _ {m}) = \det (\Phi (x _ {i}, y _ {k})) \\ (x _ {i} \in \mathrm{E}, y _ {i} \in \mathrm{F}, i = 1, \dots , m). \end{array} \right.\tag{37}
\]

定义 12. --- 给定两个 A-模 E、F 与一个半双线性型 \(\Phi\) （关于 J，在 \(E \times F\) 上），由 (37) 刻画的半双线性型 \(\Phi_{(m)}\) （关于 J，在 \((\bigwedge^{m} E) \times (\bigwedge^{m} F)\) 上）称为 \(\Phi\) 到 m 次外幂上的扩张。

在定义 12 的记号下，我们来研究与 \(\Phi_{(m)}\) 相伴的映射。由公式 (24)（第 6 小节）与 (37)，我们导出关系式

\[
\begin{array}{r l} \Phi_ {(m)} (x _ {1} \wedge \dots \wedge x _ {m}, y _ {1} \wedge \dots \wedge y _ {m}) & = \det (\langle x _ {i}, d _ {\Phi} (y _ {k}) \rangle) \\ & = \det (\langle y _ {i}, s _ {\Phi} (x _ {k}) \rangle^ {j}). \end{array}
\]

因此我们有

\[
s _ {\Phi (m)} = k _ {s} \circ (\stackrel {m} {\wedge} s _ {\Phi}), \quad d _ {\Phi_ {m}} = k _ {d} \circ (\stackrel {m} {\wedge} d _ {\Phi}),\tag{38}
\]

其中 \(k_s\) （相应地 \(k_d\) ）表示 \(\bigwedge^m F^*\) 到 \((\bigwedge^m F)^*\) 中（相应地 \(\bigwedge^m E^*\) 到 \((\bigwedge^m E)^*\) 中）的典范映射（参看第 III 章，第 8 节，第 2 小节）。

命题 10. --- 设 A 是一个配备了自同构 J 的交换域，E 与 F 是 A 上两个有限维向量空间，\(\Phi\) 是 E \(\times\) F 上的一个（关于 J 的）半双线性型。若 \(\Phi\) 非退化，则它到 m 次外幂上的扩张 \(\Phi_{(m)}\) 非退化，且 \(\Phi_{(m)}\) 的逆型是逆型 \(\hat{\Phi}\) （\(\Phi\) 的逆型）到 m 次外幂上的扩张。

事实上，由于 \(s_{\Phi}\) 与 \(d_{\Phi}\) 按假设是双射（命题 6，第 6 小节），它们的外幂也是双射（第 III 章，第 5 节，第 7 小节）。另一方面，典范映射 \(k_s\) 与 \(k_d\) 是双射（第 III 章，第 8 节，第 2 小节，定理 1）。于是，鉴于 (38)，\(s_{\Phi(m)}\) 与 \(d_{\Phi(m)}\) 是双射，这就证明了 \(\Phi_{(m)}\) 非退化（命题 6，第 6 小节）。在第二个断言中，我们已把 \(\bigwedge^m F^*\) 与 \((\bigwedge^m F)^*\) 、\(\bigwedge^m E^*\) 与 \((\bigwedge^m E)^*\) 借助映射 \(k_s\) 与 \(k_d\) （它们在此是同构（同前））隐含地等同起来。由于 \(s_{\Phi}\) 、\(d_{\Phi}\) 、\(s_{\Phi(m)}\) 与 \(d_{\Phi(m)}\) 都是双射（第 7 小节），命题中所述的逆型都存在。于是令 \(x' = x_1' \wedge \ldots \wedge x_m'\) 与 \(y' = y_1' \wedge \ldots \wedge y_m'\) （ \(x_i' \in E^*\) ， \(y_i' \in F^*\) ， \(i = 1, \ldots, m\) ）。由逆型的定义（第 7 小节）并鉴于 (38)，我们有

\[
\begin{array}{r l} \hat {\Phi} _ {(m)} (k _ {s} (y ^ {\prime}), k _ {d} (x ^ {\prime})) & = \Phi_ {(m)} (s _ {\Phi} ^ {- 1} (y _ {1} ^ {\prime}) \wedge \dots \wedge s _ {\Phi} ^ {- 1} (y _ {m} ^ {\prime}), d _ {\Phi} ^ {- 1} (x _ {1} ^ {\prime}) \wedge \dots \wedge d _ {\Phi} ^ {- 1} (x _ {m} ^ {\prime})) \\ & = \det (\Phi (s _ {\Phi} ^ {- 1} (y _ {i} ^ {\prime}), d _ {\Phi} ^ {- 1} (x _ {k} ^ {\prime}))) = \det (\hat {\Phi} (y _ {i} ^ {\prime}, x _ {k} ^ {\prime})) \end{array}
\]

由此得到我们的第二个断言。

注记. --- 设 E 是一个自由 A-模，θ 是 ∧E 到 m 阶反称张量所成的子模上的典范同构（第 III 章，第 5 节，第 6 小节，命题 6）。设 Φ 是 E 上的一个半双线性型，Φ(m) 是 Φ 到 ∧E 上的扩张，而 Θ 是 ∧E 上的这样一个半双线性型：它是 Φ 到 ⊗E 上的扩张关于 θ 的原像。由 θ 与张量的反称化的定义以及 (35)，我们有

\[
\Theta (x _ {1} \wedge \dots \wedge x _ {m}, y _ {1} \wedge \dots \wedge y _ {m}) = \sum_ {\sigma , \tau} \varepsilon_ {\sigma} \varepsilon_ {\tau} \Phi (x _ {\sigma (1)}, y _ {\tau (1)}) \dots \Phi (x _ {\sigma (m)}, y _ {\tau (m)})
\]

其中 \(\sigma\) 与 \(\tau\) 取遍对称群 \(\mathfrak{S}_{m}\) 。由行列式的计算公式与公式 (37)，这个表达式可写成

\[
\sum_ {\tau \in \mathfrak {S} _ {m}} \varepsilon_ {\tau} \det (\Phi (x _ {i}, y _ {\tau (k)})) = m! \det (\Phi (x _ {i}, y _ {k}));
\]

换句话说，我们有 \(\Theta = m!\Phi_{(m)}\) 。

\subsection{矩阵计算。}

本小节中，我们打算推广第 II 章，第 6 节中引入的矩阵演算，并用它来转述本节中已证明的某些结果。

I. --- 设 I 与 K 是两个有限指标集，H 是一个非空集合，\(M = (m_{ik})_{(i,k) \in \mathbb{I} \times \mathbb{K}}\) 是 H 上的一个矩阵（第 II 章，第 6 节，第 1 小节，定义 1）。

称为 \(M\) 的转置（记作 \({}^t M\) ）的是矩阵 \((m_{ki}')_{(k,i) \in \mathbb{K} \times \mathbf{I}}\) ，它满足 \(m_{ki}' = m_{ik} ((i, k) \in \mathbf{I} \times \mathbf{K})\) 。显然有

\[
^ {t} (^ {t} M) = M.\tag{39}
\]

这推广了第 II 章，第 6 节，第 6 小节中引入的概念。

设 H 是一个交换群（记法为加法）。以 I 与 K 为指标集的 H 上矩阵的集合具有交换群结构，因为它是 I × K 到 H 的映射的集合。这个群用加法记之。

设 H′、H″ 是两个非空集合，H 是一个交换群（记法为加法），\(f:(h',h'')\to h'h''\) 是 \(\mathbf{H}'\times\mathbf{H}''\) 到 H 中的一个映射。给定两个矩阵

\[
M ^ {\prime} = (m _ {i k} ^ {\prime}) _ {(i, k) \in \mathrm{I} \times \mathrm{K}}, \quad M ^ {\prime \prime} = (m _ {k l} ^ {\prime \prime}) _ {(k, l) \in \mathrm{K} \times \mathrm{L}}
\]

（分别在 H′ 与 H″ 上），使得 M′ 的列指标集 K 等于 M″ 的行指标集；称 M′ 与 M″ （依 f ）的积，并记作 M′M″ ，为矩阵

\[
M ^ {\prime}. M ^ {\prime \prime} = (\sum_ {k \in \mathbf {K}} m _ {i k} ^ {\prime} m _ {k l} ^ {\prime \prime}) _ {(i, l) \in \mathbf {I} \times \mathbf {L}}\tag{40}
\]

（在 H 上）。这推广了第 II 章，第 6 节，第 4 小节中引入的概念。若 \(\mathbf{H}' = \mathbf{H}'' = \mathbf{H}\) 且 H 是一个环，则除非另有明确说明，积 \(M'M''\) 将“在 H 中”计算，也就是按照映射 \((x, y) \to xy\) 计算。当 \(\mathbf{H}'\) 与 \(\mathbf{H}''\) 是交换群（记法为加法）且 \(f\) 是双线性映射时，有

\[
\left\{ \begin{array}{l} (M ^ {\prime} + M _ {1} ^ {\prime}) M ^ {\prime \prime} = M ^ {\prime} M ^ {\prime \prime} + M _ {1} ^ {\prime} M ^ {\prime \prime}, \\ M ^ {\prime} (M ^ {\prime \prime} + M _ {1} ^ {\prime \prime}) = M ^ {\prime} M ^ {\prime \prime} + M ^ {\prime} M _ {1} ^ {\prime \prime}, \end{array} \right.\tag{41}
\]

其中 \(M'\) 、\(M_1'\) 是 \(H'\) 上的矩阵，\(M''\) 、\(M_1''\) 是 \(H''\) 上的矩阵，且写出的和与积都假定有定义。设 \(M'\) 、\(M''\) 是集合 \(H'\) 、\(H''\) 上的矩阵，\(f^0\) 是 \(H'' \times H'\) 到 H 中的、由 \((h'', h') \to h' h''\) 定义的映射；则有

\[
{ } ^ { t } ( M ^ { \prime } M ^ { \prime \prime } ) = { } ^ { t } M ^ { \prime \prime } . { } ^ { t } M ^ { \prime }\tag{42}
\]

其中第一个（相应地第二个）成员中的积按照 \(f\) （相应地 \(f^{0}\) ）计算。

在 \(H' = H'' = H\) 是一个环的情形，我们重新得到第 II 章，第 6 节，第 6 小节的公式 (12)。

设 A 是一个环，J 是 A 的一个反自同构。对 A 上每个矩阵 \(M = (m_{ik})\) ，我们将用 \(M^{J}\) 表示矩阵 \((m_{ik}^{J})\) 。设 \(M_{1}\) 、\(M_{2}\) 是 A 上两个矩阵，使得 \(M_{1}M_{2}\) 有定义。由于 J 是 A 到反环 \(A^{0}\) 上的一个同构，有 \((M_{1}M_{2})^{J} = M_{1}^{J}\) . \(M_{2}^{J}\) ，其中第一个（相应地第二个）成员在 A （相应地 \(A^{0}\) ）中计算。鉴于 (42) 与 (39)，这给出

\[
(M _ {1} M _ {2}) ^ {\mathrm{J}} = ^ {t} (^ {t} M _ {2} ^ {\mathrm{J}}. ^ {t} M _ {1} ^ {\mathrm{J}})\tag{43}
\]

其中两个成员都在 A 中计算。

设 \(H_{1}\) 、\(H_{2}\) 、\(H_{3}\) 、\(H_{12}\) 、\(H_{23}\) 与 H 是交换群（记法为加法），\(f_{12}: H_{1} \times H_{2} \to H_{12}\) 、\(f_{23}: H_{2} \times H_{3} \to H_{23}\) 、\(f_{3}: H_{12} \times H_{3} \to H\) 、\(f_{1}: H_{1} \times H_{23} \to H\) 是映射，并设 \(M_{1}\) 、\(M_{2}\) 、\(M_{3}\) 分别是 \(H_{1}\) 、\(H_{2}\) 、\(H_{3}\) 上的矩阵。若有 \(f_{3}(f_{12}(x_{1}, x_{2}), x_{3}) = f_{1}(x_{1}, f_{23}(x_{2}, x_{3}))\) （对一切 \(x_{i} \in H_{i}\) ，i = 1, 2, 3），则积 \((M_{1}M_{2})M_{3}\) 与 \(M_{1}(M_{2}M_{3})\) （分别按照 \(f_{12}\) 、\(f_{3}\) 、\(f_{23}\) 与 \(f_{1}\) 计算）若有定义便相等；将把它们记作 \(M_{1}M_{2}M_{3}\) 。当 \(H_{1} = H_{2} = H_{3} = H_{12} = H_{23} = H\) 、H 是一个环且 \(f_{12}\) 、\(f_{23}\) 、\(f_{3}\) 、\(f_{1}\) 都等于映射 \((x, y) \to xy\) 时，前一条件表示后者的结合律，因此是满足的。对多于三个因子的积，将作类似的约定。

设 A、B 是两个环，\(M = (m_{ik})_{(i,k) \in I \times K}\) 与 \(M' = (m_{ik}')_{(i,k) \in I \times K}\) 是 (A, B)-双模 G 上的两个矩阵（第 1 小节）。若对每个元素在 A 中的单行矩阵 \(L = (a_i)_{i \in I}\) 与每个元素在 B 中的单列矩阵 \(C = (b_k)_{k \in K}\) ，都有 \(L.M.C = L.M'.C\) （积按照定义 G 的 (A, B)-双模结构的映射计算），则矩阵 \(M\) 与 \(M'\) 相等。事实上，若取 \(a_i = 1\) ， \(a_s = 0\) （ \(s \neq i\) ）， \(b_k = 1\) ， \(b_t = 0\) （ \(t \neq k\) ），则纯量矩阵 \(L.M.C\) 与 \(L.M'.C\) 分别等于 \(m_{ik}\) 与 \(m_{ik}'\) 。

\begin{enumerate}
\def\labelenumi{\Roman{enumi}.}
\setcounter{enumi}{1}
\tightlist
\item
  考虑一个环 A 与一个（右或左）A-模 E，它具有有限基 \((e_{i})_{i\in I}\) 。对 E 的每个元素 x，由分量 \(x_{i}\) （ \(i\in I\) ，即 x 关于 \((e_{i})\) 的分量）构成的单列矩阵称为 x 关于基 \((e_{i})\) 的矩阵，记作 \(M(x)\) 或 x（参看第 II 章，第 6 节，第 4 小节）；在计算中，为记起指标 i 是行指标，给它添上一个只能取一个值的列指标，并用 \((x_{i0})\) 来记矩阵 \(M(x)\) ，将是方便的。
\end{enumerate}

现在考虑两个（左或右）A-模 E 与 F，分别具有有限基 \((e_{i})_{i\in I}\) 与 \((f_{k})_{k\in K}\) ；设 \((f_{k}^{*})\) 是 \(F^{*}\) 的与 \((f_{k})\) 对偶的基。我们将在下列四种情形中，定义 E 到 F 中的映射 u 关于这些基的矩阵：

(D) E 与 F 是右模，且 u 是 A-线性的；

(G) E 与 F 是左模，且 u 是 A-线性的；

(GD) E 是左模，F 是右模，A 配备了一个反自同构 J，u 是 Z-线性的且满足 \(u(ax) = u(x)a^{J}\) （ \(a \in A, x \in E\) ）（换句话说，u 是 \(E^{J}\) 到 F 中的 A-线性映射（第 2 小节，定义 5））。

(DG) E 是右 A-模，F 是左 A-模，A 配备了一个反自同构 J，u 是 Z-线性的且满足 \(u(xa) = a^{J}u(x)\) （ \(x \in E, a \in A\) ）（换句话说，u 是 \(E^{J}\) 到 F 中的一个 A-线性映射）。

在这四种情形的每一种中，映射 u 的矩阵按定义是使得下式成立的矩阵 \((u_{ki})_{(k,i)\in\mathbb{K}\times\mathbb{I}}\) ：

\[
u _ {k i} = \left\langle u (e _ {i}), f _ {k} ^ {*} \right\rangle .\tag{44}
\]

在情形 (D)，这个定义与第 II 章，第 6 节，第 3 小节给出的定义一致。在这些条件下，矩阵 \(M(u(x))\) （E 的元素 \(x\) 的像的矩阵）由下列公式给出：

\[
M (u (x)) = M (u). M (x)\tag{45 D}
\]

\[
{ } ^ { t } M ( u ( x ) ) = { } ^ { t } M ( x ) . { } ^ { t } M ( u )\tag{45 G}
\]

\[
M (u (x)) = M (u). M (x) ^ {\mathrm{J}}\tag{45 GD}
\]

\[
{ } ^ { t } M ( u ( x ) ) = { } ^ { t } M ( x ) ^ { j } . { } ^ { t } M ( u ) .\tag{45 DG}
\]

例如验证 (45 DG)，其余验证类似且略更容易。设 \(x = \sum_{i} e_i x_{io}, u(x) = \sum_{k} y_{ko} f_k\) ；我们有 \(u(x) = u(\sum_{i} e_i x_{io}) = \sum_{i} x_{io}^{\mathbf{J}} u(e_i) = \sum_{i,k} x_{io}^{\mathbf{J}} u_{ki} f_k\) ；由此得 \(y_{ko} = \sum_{i} x_{io}^{\mathbf{J}} u_{ki}\) ；为使两个指标 \(i\) 彼此靠近，考虑转置矩阵 \(^t M(x) = (x_{oi}')\) （其中 \(x_{oi}' = x_{io}\) ）与 \(^t M(u) = (u_{ik}')\) （其中 \(u_{ik}' = u_{ki}\) ）；于是有 \(y_{ko} = \sum_{i} x_{oi}' u_{ik}'\) ；由于右端是以 \(k\) 为指标的元素（属于单行矩阵 \(^t M(x)^{\mathbf{J}}\) . \(^t M(u)\) ），公式 (45 DG) 得到验证。

注记. -- 1) 当 A 交换时，借助公式 \(t(M'M'') = 'M'''.'M'\) （参看 (42)），(45 G) 化为 (45 D)，(45 DG) 化为 (45 GD)，其中两端的成员这里都在 A 中计算。2) 设 E、F、G 是三个具有有限基的左模，\(u: \mathrm{E} \to \mathrm{F}\) 、\(\nu: \mathrm{F} \to \mathrm{G}\) 是 A-线性映射。由 (45 G) 得出

\[
{ } ^ { t } M ( \nu \circ u ) = { } ^ { t } M ( u ) . { } ^ { t } M ( \nu ) .\tag{46}
\]

事实上，对任意 \(x \in E\) ，我们有

\[
\begin{array}{r l} ^ {t} M (x). ^ {t} M (\nu \circ u) & = ^ {t} M (\nu (u (x))) = ^ {t} M (u (x)). ^ {t} M (\nu) \\ & = ^ {t} M (x). ^ {t} M (u). ^ {t} M (\nu), \end{array}
\]

由此得 (46)。

回顾在右模的情形，我们有

\[
M (\nu \circ u) = M (\nu) M (u).
\]

\begin{enumerate}
\def\labelenumi{\Roman{enumi}.}
\setcounter{enumi}{2}
\tightlist
\item
  从现在起，我们用 A 表示一个环，用 B 表示一个环（相应地一个配备了反自同构 J 的环，对此我们令 \(J' = J^{-1}\) ），用 E 表示具有有限基 \((e_i)_{i \in I}\) 的一个左 A-模，用 F 表示具有有限基 \((f_k)_{k \in K}\) 的一个右（相应地左）B-模。我们用 \((e_i^*)\) 与 \((f_k^*)\) 表示 \(E^*\) 与 \(F^*\) 的对偶基。除非另有明确说明，所考虑的矩阵都关于这些基取定。
\end{enumerate}

设 G 是一个 (A, B)-双模（第 1 小节），Φ 是 E × F 到 G 中的一个双线性（相应地关于 J 的右半双线性）映射，\(R = (\Phi(e_i, f_k))\) 是 Φ 的矩阵。于是，对 \(x \in E\) 与 \(y \in F\) ，第 1 小节的公式 (6)（相应地第 2 小节的 (8)），鉴于上述约定，写成

\[
\Phi (x, y) = ^ {t} M (x). R. M (y) \quad (\text { resp. } \Phi (x, y) = ^ {t} M (x). R. M (y) ^ {\mathrm{J}}),\tag{47}
\]

其中积按照定义 (A, B)-双模 G 的结构的映射计算；特别地，若 A = B = G（这时 Φ 是一个型），则积在 A 中计算。

设 E′ 是具有有限基 \((e_{s}^{\prime})_{s\in S}\) 的一个左 A-模，F′ 是具有有限基 \((f_{t}^{\prime})_{t\in T}\) 的一个右（相应地左）A-模，\(u:E\to E^{\prime}\) 与 \(\nu:F\to F^{\prime}\) 是 A-线性映射，\(\Phi^{\prime}\) 是 \(E^{\prime}\times F^{\prime}\) 到 G 中的一个双线性（相应地关于 J 的右半双线性）映射。用 \(\Phi\) 表示 \(\Phi^{\prime}\) （关于 u 与 \(\nu\) ）的原像，用 U、V、R、\(R^{\prime}\) 表示 u、\(\nu\) 、\(\Phi\) 、\(\Phi^{\prime}\) 关于所考虑的基的矩阵。于是有

\[
R = ^ {t} U. R ^ {\prime}. V \quad (\text { resp. } R = ^ {t} U. R ^ {\prime}. V ^ {J}),\tag{48}
\]

其中积如 (47) 中那样计算。事实上，无论 \(x \in \mathbf{E}\) 与 \(y \in \mathbf{F}\) 如何，按定义有 \(\Phi(x, y) = \Phi'(u(x), v(y))\) ，于是由 (47)，

\[
\begin{array}{c} ^ {t} M (x). R. M (y) = ^ {t} M (u (x)). R ^ {\prime}. M (\nu (y)) \\ (\text {resp.} ^ {t} M (x). R. M (y) ^ {\mathrm{J}} = ^ {t} M (u (x)). R ^ {\prime}. M (\nu (y)) ^ {\mathrm{J}}); \end{array}
\]

由 (45 G) 与 (45 D)（相应地 (45 G)）以及 (43)，导出

\[
\begin{array}{r l} ^ {t} M (x). R. M (y) & = ^ {t} M (x). ^ {t} U. R ^ {\prime}. V. M (y) \\ (\text {resp.} ^ {t} M (x). R. M (y) ^ {J} & = ^ {t} M (x). ^ {t} U. R ^ {\prime}. ^ {t} (^ {t} M (y). ^ {t} V) ^ {J} \\ & = ^ {t} M (x). ^ {t} U. R ^ {\prime}. V ^ {J}. M (y) ^ {J}); \end{array}
\]

这就证明了我们的断言。

\begin{enumerate}
\def\labelenumi{\Roman{enumi}.}
\setcounter{enumi}{3}
\tightlist
\item
  --- 这里假设环 A 与 B 相等，并用 \(\Phi\) 表示 E × F 上的一个双线性（相应地关于 J 的右半双线性）型，用 \(R\) 表示其矩阵。我们来计算映射 \(s_{\Phi}\) 与 \(d_{\Phi}\) （与 \(\Phi\) 相伴者）的矩阵，为简单起见将把它们记作 \(s\) 与 \(d\) 。由于由第 6 小节的 (23)（相应地由第 6 小节的 (24)）有 \(\Phi(x, y) = \langle y, s(x) \rangle = \langle x, d(y) \rangle\) （相应地 \(\Phi(x, y) = \langle x, d(y) \rangle = \langle y, s(x) \rangle^{J}\) ），便有 \(\Phi(e_i, f_k) = \langle f_k, s(e_i) \rangle = \langle e_i, d(f_k) \rangle\) （相应地 \(\Phi(e_i, f_k) = \langle e_i, d(f_k) \rangle = \langle f_k, s(e_i) \rangle^{J}\) ），于是，由 (44) 并由于 ( \(e_i\) ) 是 ( \(e_i^*\) ) 的对偶基、( \(f_k\) ) 是 ( \(f_k^*\) ) 的对偶基：
\end{enumerate}

\[
M (d) = R,   M (s) = ^ {t} R \quad (\text { resp. }   M (d) = R,   M (s) = ^ {t} R ^ {J ^ {\prime}}).\tag{49}
\]

注记. --- 1) 当 A 是一个体时，线性映射 s 与 d 有相同的秩。我们在这里看到，它们的矩阵 \(M(s)\) 与 \(M(d)\) 有相同的秩；事实上，A 上的矩阵与其转置有相同的秩（第 II 章，第 6 节，第 7 小节，命题 3），且当 \(\Phi\) 是半双线性时，\(R\) 在 A 上的秩与 \(^tR\) 在 A \(^0\) 上的秩相等（同处）以及 J′ 是 A \(^0\) 到 A 上的同构这一事实，蕴涵 \(R\) 与 \(^tR^{J'}\) 在 A 上的秩相等。

\begin{enumerate}
\def\labelenumi{\arabic{enumi})}
\setcounter{enumi}{1}
\tightlist
\item
  若 M 与 N 是具有有限基 \((m_{i})\) 与 \((n_{k})\) 的右 A-模，\(\Phi\) 是 \(M \times N\) 上的一个（关于 J 的）左半双线性型（第 6 小节，注记），s 与 d 是它的相伴映射，\(R = (\Phi(m_{i}, n_{k}))\) 是它的矩阵，则第 6 小节的公式 (26) 表明有
\end{enumerate}

\[
M (d) = R ^ {J ^ {\prime}}, \quad M (s) = ^ {t} R.
\]

现在设与 \(\Phi\) 相伴的映射 s 与 d 是双射，而来计算矩阵 \(\hat{R}\) （\(\Phi\) 的逆型的矩阵，第 7 小节）。当 \(\Phi\) 是双线性型时，\(\Phi\) 是 \(\hat{\Phi}\) 关于线性映射 \(s: \mathrm{E} \to \mathrm{F}^*\) 与 \(d: \mathrm{F} \to \mathrm{E}^*\) 的原像；于是，鉴于 (48) 与 (49)，有 \(R = R\) . \(\hat{R}\) . \(R\) ，从而，由于 \(R\) 可逆（ \(d\) 是双射），\(\hat{R} = R^{-1}\) 。这个公式可推广到 \(\Phi\) 是半双线性型的情形，因为，若把 \(\Phi\) 看作 \(\mathrm{E} \times \mathrm{F}^{\mathrm{j}}\) 上的双线性型，并把 \((\mathrm{F}^{\mathrm{j}})^{*}\) 与 \((\mathrm{F}^{*})^{\mathrm{j}}\) 等同起来（参看第 9 小节），则这个双线性型的逆型与 \(\hat{\Phi}\) （看作 \((\mathrm{F}^{*})^{\mathrm{j}} \times \mathrm{E}^{*}\) 上的双线性型）一致。在两种情形下，\(\Phi\) 的逆型的矩阵都是 \(\Phi\) 的矩阵的逆。

最后，设 E′ 是一个左 A-模，F′ 是一个右（相应地左）A-模，两者都具有有限基 \((e_{s}^{\prime})\) 与 \((f_{t}^{\prime})\) ；设 \(\Phi^{\prime}\) 是 \(E^{\prime} \times F^{\prime}\) 上的一个双线性型（相应地关于 J 的半双线性型），\(R^{\prime}\) 是它的矩阵。设 \(s_{\Phi}\) 与 \(d_{\Phi}\) 是双射。设 \(u : E \to E^{\prime}\) 与 \(\nu : F \to F^{\prime}\) 是线性映射，\(u^{*} : F^{\prime} \to F\) 与 \(\upsilon^{*} : E^{\prime} \to E\) 是它们的伴随（第 8 小节，命题 7）；用 U、V、\(U^{*}\) 、\(V^{*}\) 表示 u、\(\upsilon\) 、\(u^{*}\) 、\(\upsilon^{*}\) 关于给定基的矩阵。于是有

\[
U ^ {*} = R ^ {- 1}. ^ {t} U. R ^ {\prime}, \quad {} ^ {t} V ^ {*} = R ^ {\prime}. V. R ^ {- 1}\tag{50}
\]

\[
(\text { resp. } U ^ {* \mathrm{J}} = R ^ {- 1}. ^ {t} U. R ^ {\prime}, ^ {t} V ^ {*} = R ^ {\prime}. V ^ {\mathrm{J}}. R ^ {- 1}).
\]

事实上，无论 \(x \in E\) 与 \(y \in F'\) 如何，有 \(\Phi'(u(x), y) = \Phi(x, u^*(y))\) （第 8 小节，定义 10）。于是，当 \(\Phi\) 是双线性型时，鉴于 (47)，有 \(^t M(u(x)). R'. M(y) = ^t M(x). R. M(u^*(y))\) ；这给出，鉴于 (45 G) 与 (45 D)，\(^t M(x). ^t U. R'. M(y) = ^t M(x). R. U^*. M(y)\) ，从而 \(^t U. R' = R. U^*\) 以及第一个所宣布的公式，因为 \(d\) 是双射，\(R\) 可逆。当 \(\Phi\) 是半双线性型时，(47) 与 (45 G) 给出 \(^t M(x). ^t U. R'. M(y)^J = ^t M(x). R. (^{t}M(y). ^{t}U^*)^{J}\) ；而由 (43)，我们有 \((^{t}M(y). ^{t}U^*)^{J} = ^{t}(^{u}U^{*J}. ^{u}M(y)^{J})\) ，从而 \((^{t}M(y). ^{t}U^*)^{J} = U^{*J}. M(y)^{J}\) ；因此导出 \(^t M(x). ^{t}U. R'. M(y)^{J} = ^{t}M(x). R. U^{*J}. M(y)^{J}\) ，从而 \(^t U. R' = R. U^{*J}\) ，以及 \(U^{*J} = R^{-1}. ^{t}U. R'\) 。关于 \(V^*\) 的公式的验证是类似的。

习题. --- 1) 设 A 是一个交换域，E 是 A 上的一个向量空间，具有可数无穷基 \((e_n)_{n \geqslant 1}\) 。如下定义 E 上的一个双线性型 \(\Phi\) ：令 \(\Phi(e_{i+1}, e_i) = 1\) （对 \(i \geqslant 1\) ），\(\Phi(e_k, e_j) = 0\) （对 \(k \neq j + 1\) 与 \(j \geqslant 1\) ）。证明线性映射 \(d_\Phi\) （与 \(\Phi\) 右相伴者）是单射，但线性映射 \(s_\Phi\) （与 \(\Phi\) 左相伴者）不是单射。

\begin{enumerate}
\def\labelenumi{\arabic{enumi})}
\setcounter{enumi}{1}
\item
  设 E 是 Z 与 Z/(2) 的 Z-模直和，E* 是它的对偶（同构于 Z）。证明双线性型 \((x, x') \to \langle x, x' \rangle\) （在 \(E \times E^*\) 上）使得右相伴线性映射是单射，而左相伴线性映射不是。
\item
  给出一个双线性型 \(\Phi\) 的例子，它定义在两个向量空间的积 \(\mathbf{E} \times \mathbf{F}\) 上，使得 \(d_{\Phi}\) 是双射，而 \(s_{\Phi}\) 是单射但非双射（取 \(\mathbf{E}\) 为无限维，且 F 等于 \(\mathbf{E}^*\) ，即 \(\mathbf{E}\) 的对偶；参看第 II 章，第 5 节，习题 3）。
\item
  设 A 是一个配备了反自同构 J 的环，E 是一个左 A-模，G 是一个 (A, A)-双模，\(\Phi\) 是 \(\mathrm{E} \times \mathrm{E}\) 到 G 中的一个（关于 J 的右）半双线性映射。证明恒等式（其中 \(\mathrm{Q}(x) = \Phi(x, x)\) ）：
\end{enumerate}

\[
\begin{array}{r l} & {2 \Phi (x, y) (\mu^ {\mathrm{J}} \lambda^ {\mathrm{J}} - \lambda^ {\mathrm{J}} \mu^ {\mathrm{J}}) = \mathrm{Q} (x - \mu \lambda y) - \mathrm{Q} (x + \mu \lambda y) + \mu \mathrm{Q} (x + \lambda y)} \\ & {- \mu \mathrm{Q} (x - \lambda y) + \mathrm{Q} (x + \mu y) \lambda^ {\mathrm{J}} - \mathrm{Q} (x - \mu y) \lambda^ {\mathrm{J}} + \mu \mathrm{Q} (x - y) \lambda^ {\mathrm{J}} - \mu \mathrm{Q} (x + y) \lambda^ {\mathrm{J}}.} \end{array}
\]

\begin{enumerate}
\def\labelenumi{\arabic{enumi})}
\setcounter{enumi}{4}
\tightlist
\item
  设 K 是一个特征为 2 的交换域，A 是 K 的一个可分二次扩张；我们有 \(\mathrm{A} = \mathrm{K}(\theta)\) ，其中 \(\theta\) 是 \(\mathrm{X}^2 +\mathrm{X} + \beta\) （K{[}X{]} 的一个不可约多项式）的根，而 A 的异于恒同的 K-自同构 J 满足 \(\theta^{j} = \theta +1\) （第 V 章，第 11 节，习题 8）。证明：若 E 与 G 是 A 上的向量空间，\(\Phi\) 是 \(\mathrm{E}\times \mathrm{E}\) 到 G 中的一个（关于 J 的）半双线性映射，则令 \(\mathrm{Q}(x) = \Phi (x,x)\) 后，有
\end{enumerate}

\[
\Phi (x, y) = \mathrm{Q} (\theta x + y) - \beta \mathrm{Q} (x) - \mathrm{Q} (y) - (\theta + 1) (\mathrm{Q} (x + y) - \mathrm{Q} (x) - \mathrm{Q} (y)).
\]

\begin{enumerate}
\def\labelenumi{\arabic{enumi})}
\setcounter{enumi}{5}
\tightlist
\item
  设 A 是一个体，E 是 A 上的一个向量空间，\(\Phi\) 是 E 上的一个半双线性型，\(u\) 是 E 的一个自同态。
\end{enumerate}

\begin{enumerate}
\def\labelenumi{\alph{enumi})}
\item
  存在且仅存在 E 的一个自同态 \(u^*\) ，使得 \(\Phi(u(x), y) = \Phi(x, u^*(y))\) （对 E 中的 \(x, y\) ），必须且只需 \(d_\Phi\) 是单射且 \(^t u(d_\Phi(E)) \subset d_\Phi(E)\) 。
\item
  给出一个例子，其中 E 是无限维且 \(d_{\Phi}\) 是单射，但 \(u(d_{\Phi}(E))\) 不包含于 \(d_{\Phi}(E)\) 。
\end{enumerate}

\begin{enumerate}
\def\labelenumi{\arabic{enumi})}
\setcounter{enumi}{6}
\tightlist
\item
  设 E、\(E_{1}\) 是两个 A-模，\(\Phi\) （相应地 \(\Phi_{1}\) ）是 E （相应地 \(E_{1}\) ）上的一个半双线性型。假设 \(\Phi_{1}\) 非退化，且存在一个元素 \(\alpha\in A\) 与 E 到 \(E_{1}\) 上的一个双射 u，使得 \(\Phi_{1}(u(x), u(y))=\Phi(x,y)\alpha\) （对 E 中一切的 x, y）。证明：\(1^{\circ}\Phi\) 非退化；\(2^{\circ}u\) 是线性的；\(3^{\circ}\) 若 \(E_{1}\) 是忠实 A-模，则 E 也是，且 \(\alpha\) 不是 A 中的右零因子；\(4^{\circ}\) 若 \(\Phi_{1}\) 在 A 中取的值都不是左零因子，则 \(\Phi\) 亦然。
\end{enumerate}

¶ 8) 设 A 是一个体，E₁、E₂ 是 A 上两个非零向量空间，Φ₁ （相应地 Φ₂）是 E₁ （相应地 E₂）上的一个非退化半双线性型（关于 A 的一个反自同构 J₁ （相应地 J₂））。设 u 是 E₁ 到 E₂ 上的一个线性映射，使得关系 Φ₁(x, y) = 0 蕴涵 Φ₂(u(x), u(y)) = 0。

\begin{enumerate}
\def\labelenumi{\alph{enumi})}
\item
  证明 \(u\) 是 \(\mathbf{E}_{1}\) 到 \(\mathbf{E}_{2}\) 上的一个双射。（若 \(u(0)\) 不化为 0，证明 \(\mathbf{E}_{1}\) 中会存在两个向量 \(a\) 、\(b\) ，使得 \(u(a) \neq 0\) 、\(u(b) = 0\) 且 \(\Phi_{1}(a, b) \neq 0\) ；若 H 是由 \(x \in \mathbf{E}_{1}\) （使 \(\Phi_{1}(a, x) = 0\) 者）所成的超平面，注意将有 \(u(\mathrm{H}) = \mathrm{E}_{2}\) 。）
\item
  证明：若 dim \(E_{1} \geqslant 2\) ，则存在 \(\alpha \in A\) ，使得有 \(\Phi_{2}(u(x), u(y)) = \Phi_{1}(x, y) \alpha\) （对 \(E_{1}\) 中一切的 x, y）。（对每个 \(y \in E_{1}\) ，证明存在一个元素 \(m(y) \in A\) ，使得 \(\Phi_{2}(u(x), u(y)) = \Phi_{1}(x, y)m(y)\) （对一切 \(x \in E_{1}\) ），且若 y 与 \(y'\) 在 \(E_{1}\) 中线性无关，则有 \(m(y + y') = m(y) = m(y')\) 。）
\end{enumerate}

\begin{enumerate}
\def\labelenumi{\arabic{enumi})}
\setcounter{enumi}{8}
\tightlist
\item
  设 A 是一个体，E、F 是 A 上两个左向量空间，\(\Phi\) 是 \(E \times F\) 上的一个非退化半双线性型（关于 A 的一个反自同构 J）。
\end{enumerate}

\begin{enumerate}
\def\labelenumi{\alph{enumi})}
\item
  设 M 是 E 的一个子空间，N 是 F 的一个子空间，使得 \(N \supset M^{0}\) 且 \(M \supset N^{0}\) 。证明：若空间 \(N/M^{0}\) 、\(M/N^{0}\) 之一是有限维的，则另一个也是，且这两个空间的维数相等。
\item
  设 M、\(M'\) 是 E 的两个子空间，使得 \(M^{00} = M\) 且 \(M'\) 是有限维的；证明我们有 \((M \cap M')^{0} = M^{0} + M'^{0}\) 与 \((M + M')^{00} = M + M'\) 。（把 a) 应用于子空间 \(M'\) 与 \(M^{0} + M'^{0}\) ，证明 \(\dim(M \cap M') = \text{codim}(M^{0} + M'^{0})\) ；把 a) 应用于子空间 \(M + M'\) 与 \(M^{0}\) ，证明
\end{enumerate}

\[
\dim ((M + M ^ {\prime}) ^ {0 0} / M) = \dim ((M + M ^ {\prime}) / M).
\]

\begin{enumerate}
\def\labelenumi{\alph{enumi})}
\setcounter{enumi}{2}
\item
  若 E = F 且 M 是 E 的一个子空间，使得 \(E = M^{0} + M^{00}\) ，证明 E 是 \(M^{0}\) 与 \(M^{00}\) 的直和。
\item
  设 E 是交换域 A 上的一个向量空间，具有可数无穷基 \((e_{n})_{n\geqslant0}\) ，并设 \(\Phi\) 是 E 上的对称双线性型，使得 \(\Phi(e_{n}, e_{n}) = 1\) （对一切 n），\(\Phi(e_{i}, e_{j}) = 0\) （对 \(i \geqslant 1\) 、\(j \geqslant 1\) 且 \(i \neq j\) ），且 \(\Phi(e_{0}, e_{n}) = 1\) （对一切 \(n \geqslant 1\) ）。证明 \(\Phi\) 非退化。设 M （相应地 N）是 E 的由诸 \(e_{2k}\) （相应地诸 \(e_{2k-1}\) ）（\(k \geqslant 1\) ）生成的子空间，并设 H = M + N，它是 E 的一个超平面。证明我们有 \(M^{0} = N\) 、\(N^{0} = M\) 、\(H^{00} = E \neq H\) 、\((M \cap N)^{0} \neq M^{0} + N^{0}\) 与 \((M + N)^{00} \neq M + N\) ，尽管 \(M^{00} = M\) 、\(N^{00} = N\) ；若 L 是由 \(e_{0}\) 与 \(e_{1}\) 生成的 2 维子空间，则有 \((L \cap H)^{0} \neq L^{0} + H^{0}\) 。
\end{enumerate}

¶ 10) 设 E、E′ 分别是体 A、A′ 上两个左向量空间，维数 ≥3；设 ℱ(E) （相应地 ℱ(E′)）是 E （相应地 E′）的有限维子空间所构成的（关于包含关系的）格序集。

\begin{enumerate}
\def\labelenumi{\alph{enumi})}
\item
  设 \(p\) 是 \(\mathfrak{F}(\mathrm{E})\) 到 \(\mathfrak{F}(\mathrm{E}')\) 中的一个映射，使得对每个 \(\mathbf{M} \in \mathfrak{F}(\mathrm{E})\) 有 \(\dim p(\mathbf{M}) = \dim \mathbf{M}\) ，且对 \((\mathbf{M}, \mathbf{N})\) （\(\mathfrak{F}(\mathrm{E})\) 的元素的对）有 \(p(\mathbf{M} + \mathbf{N}) = p(\mathbf{M}) + p(\mathbf{N})\) 。证明 \(p\) 是单射；若 \(p\) 是双射，则存在一个双射半线性映射 \(u\) （从 \(\mathrm{E}\) 到 \(\mathrm{E}'\) 中），使得有 \(u(\mathbf{M}) = p(\mathbf{M})\) （对一切 \(\mathbf{M} \in \mathfrak{F}(\mathrm{E})\) ）（利用第 II 章，第 2 版，附录 III，习题 10）。
\item
  给出一个例子，其中 \(A' = A\) 是交换的，\(E' = E\) 是有限维的，且存在一个 \(\mathfrak{F}(E)\) 到其自身的映射 p，使得 \(\dim p(M) = \dim M\) 、\(p(M + N) = p(M) + p(N)\) 、\(p(M \cap N) = p(M) \cap p(N)\) ，但不存在 E 到其自身的单射半线性映射 u，使得 \(u(M) = p(M)\) （对 \(M \in \mathfrak{F}(E)\) ）。（考虑存在 A 的一个有限次数且同构于 A 的扩体 \(A''\) 的情形，例如 \(\mathbf{A} = \mathbf{F}_{p}(\mathbf{X})\) ，其中 \(p\) 是素数；这时 \(\mathrm{E}'' = \mathrm{A}'' \otimes_{\mathbf{A}} \mathrm{E}\) 是 \(\mathrm{A}''\) 上的向量空间，其维数与 \(\mathrm{E}\) 在 \(\mathrm{A}\) 上的维数相同；考虑映射 \(\mathrm{M} \to \mathrm{A}'' \otimes_{\mathbf{A}} \mathrm{M}\) 。
\item
  我们假设 \(A' = A\) 。设 \(\omega\) 是 \(\mathfrak{F}(E)\) 到 \(\mathfrak{F}'(E')\) （\(E'\) 的有限余维数子空间所成之集）上的一个双射，使得 \(\text{codim } \omega(M) = \dim M\) ，且 \(\omega(M + N) = \omega(M) \cap \omega(N)\) 。证明存在一个半双线性型 \(\Phi\) （在 \(E \times E'\) 上，右非退化），且使得 \(\omega(M) = M^0\) （对一切 \(M \in \mathfrak{F}(E)\) ）。（利用第 II 章，第 4 节，第 6 小节，定理 1）。
\end{enumerate}

¶ 11) a) 设 A 是一个左、右阿廷环（第 VIII 章，第 2 节，第 3 小节）。证明下列条件等价：1° 每个左理想（相应地右理想）≠ A 的右零化子（相应地左零化子）不化为 0；2° 每个单左 A-模（相应地右 A-模）的对偶不化为 0；3° 每个有限生成左 A-模（相应地右 A-模）的对偶不化为 0。称环 A 满足条件 (N \(_{s}\) ) （相应地 (N \(_{a}\) )），如果它对左 A-模（相应地右 A-模）满足这些条件。

\begin{enumerate}
\def\labelenumi{\alph{enumi})}
\setcounter{enumi}{1}
\item
  设 A 是一个满足条件 (N \(_{a}\) ) 的左、右阿廷环，E （相应地 F）是 A 上具有可数基的自由左 A-模（相应地右 A-模），Φ 是 E × F 上的一个双线性型，使得 s \(_{\Phi}\) 是单射。设 M 是 E 的一个自由子模；证明存在 F 的一个自由子模 N 与 M （相应地 N）的基 (e \(_{n}\) ) （相应地 (f \(_{n}\) )），使得 Φ(e \(_{i}\) , f \(_{j}\) ) = δ \(_{ij}\) ；此外，对 e \(_{1}\) 可取 M 的任意自由元素。（注意若 x 是 M 的自由元素，则 F 在映射 y → Φ(x, y) 下的像是整个环 A；然后用归纳法构造基 (e \(_{n}\) ) 与 (f \(_{n}\) )。）由此导出：若 M 的基 (e \(_{n}\) ) 是有限的，则 E 是 M 与 N \(^{0}\) 的直和，且有 M \(^{00}\) = M。
\item
  我们保持 \(b\) 的假设，并再设 A 满足条件 \((\mathrm{N}_{s})\) 且 \(d_{\Phi}\) 是单射。证明这时 E 中存在一个基 \((e_{n})\) 且 F 中存在一个基 \((f_{n})\) ，使得 \(\Phi(e_{i}, f_{j}) = \delta_{ij}\) 。（用 \(b\) ，以归纳法交替地确定 \(e_{n}\) 与 \(f_{n}\) 。）
\end{enumerate}

¶ 12) 设 E 是一个可数型实希尔伯特空间，\(\Phi(x, y)\) 是 E 中的标量积。证明：不存在 R 上的向量空间 E 的两个代数基 \((e_{\lambda})\) 、\((f_{\mu})\) 组成的组，使得对每一对指标有 \(\Phi(e_{\lambda}, f_{\mu}) = \delta_{\lambda y}\) 。（首先注意这些基的指标集将具有连续统的势（《拓扑向量空间》，第 II 章，第 3 节，习题 15）；然后考虑 E 的一个（可数）正交基 \((a_n)\) ，并注意由诸 \(a_n\) 生成的子空间包含于 \((e_{\lambda})\) 的一个可数子族所生成的子空间中。
